\documentclass[3p,final,10pt]{elsarticle}
\journal{Results in Engineering}

% ---------- packages (kept to a widely-available core) ----------
\usepackage[utf8]{inputenc}
\usepackage[T1]{fontenc}
\usepackage{lmodern}
\usepackage{amsmath,amssymb,mathtools}
\usepackage{booktabs}
\usepackage{multirow}
\usepackage{array}
\usepackage{xcolor}
\usepackage{caption}
\usepackage{enumitem}
\usepackage[hidelinks]{hyperref}
\usepackage{xurl}
\usepackage{graphicx}
\usepackage{tikz}
\usetikzlibrary{positioning,fit,backgrounds,calc,arrows.meta}
\usepackage{algorithm}
\usepackage{algpseudocode}
\usepackage{hyphenat}
\usepackage{lineno}
\modulolinenumbers[5]
\usepackage{float}
\usepackage{placeins}
% keep every figure/table/algorithm exactly where it is written, instead of
% letting LaTeX's default float rules defer them (they were all piling up
% at the end of the document otherwise)
\let\oldsection\section
\renewcommand{\section}{\FloatBarrier\oldsection}
\let\oldsubsection\subsection
\renewcommand{\subsection}{\FloatBarrier\oldsubsection}

\definecolor{probelane}{RGB}{234,238,253}
\definecolor{causallane}{RGB}{253,236,227}
\definecolor{boxline}{RGB}{104,112,138}
\definecolor{inkcolor}{RGB}{35,45,90}
\definecolor{paperbg}{RGB}{243,239,230}
\definecolor{probeaccent}{RGB}{67,97,238}
\definecolor{causalaccent}{RGB}{243,114,44}
\definecolor{stepnum}{RGB}{67,97,238}
\definecolor{diagshadow}{RGB}{0,0,0}
\definecolor{stagebg}{RGB}{246,246,244}
\definecolor{stageborder}{RGB}{206,206,200}
\definecolor{stagetext}{RGB}{60,60,58}

\captionsetup{font=small,labelfont=bf}

\begin{document}

\begin{frontmatter}

\title{Localising subject-identity and acquisition leakage in Parkinson's
handwriting classifiers: background suppression is verified but
insufficient}

\author[1]{Jesu Arockiaraj J}
\ead{jesuarockiaraj.j@vit.ac.in}
\author[1]{Arulkumar V\corref{cor1}}
\ead{arulkumar.v@vit.ac.in}
\cortext[cor1]{Corresponding author.}

\affiliation[1]{organization={School of Computer Science and Engineering, Vellore Institute of Technology},
            city={Vellore},
            state={Tamil Nadu},
            country={India}}

% ======================================================================
\begin{abstract}
\noindent
\textbf{Background.} Deep-learning Parkinson's disease (PD) handwriting
classifiers routinely report 95--99\% accuracy under image-level
train/test splitting. We replicate this evaluation gap under one matched
protocol and separate two phenomena: a small subject-identity leakage
effect (image-level vs.\ subject-grouped splitting) and a much larger
cross-cohort acquisition/domain-shift effect (within-cohort vs.\ external
validation). Unestablished is where the domain-shift signal lives background-confined signal is removable by preprocessing trace-intrinsic
signal is not.
\textbf{Methods.} We probed frozen ImageNet-pretrained EfficientNet-B0
features under three masking modes (full/background-only/trace-only) to
localise subject re-identification and cohort decodability, built a
background/template suppression pipeline, confirmed it removed the
targeted signal, and tested whether honest-protocol AUC improved via a
paired, subject-level bootstrap. A rendering-invariant, geometry-only
reconstruction ruled out a colour confound.
\textbf{Results.} Background alone perfectly separated cohorts, surviving
unchanged under grayscale and geometry-only renderings (cohort AUC
$1.000\to1.000\to0.997$), ruling out colour. Suppression removed
background-confined signal (re-identification and cohort decodability
collapsed to chance) but produced no improvement in honest-protocol AUC in
any cross-cohort or within-cohort comparison (all 95\% CIs included
zero). Given this design's minimum detectable effect
($\approx0.13$--$0.37$ AUC), this is evidence of absence, not evidence of
failure. One unreplicated Deep CORAL test likewise showed no improvement.
\textbf{Conclusion.} The identity and cohort signal this study localises is
geometric and trace-intrinsic, not background-confined or colour-related
suppressing the background/template channel specifically does not improve
honest-protocol accuracy, so that one class of preprocessing-only
de-leaking intervention is insufficient on this evidence. Cohort
decodability from the trace is correlational, not a causal attribution of
the cross-cohort AUC drop to any single mechanism population, disease-stage,
or protocol differences between cohorts remain equally compatible
explanations and are not ruled out here. We
recommend pairing any de-leaking intervention with a signal-localisation
probe run before and after, including a rendering-invariance control,
released as a reusable protocol.
\end{abstract}

\begin{keyword}
Parkinson's disease \sep handwriting \sep data leakage \sep subject re-identification \sep mechanistic evaluation \sep domain shift
\end{keyword}

\end{frontmatter}

\linenumbers

% ======================================================================
\section{Introduction}

Deep learning classifiers trained on hand-drawn spirals, meanders and waves
routinely report 95--99\% accuracy for Parkinson's disease (PD) screening
\cite{gazda2022,kamran2021}. This is almost always measured by pooling all
drawings and splitting them at the image level into training and test sets
because each participant contributes several drawings, this procedure lets
drawings from the same person land on both sides of the split, so a network
can succeed by recognising the individual rather than the disease.

As a first step, we replicate this evaluation gap ourselves, under a single
protocol held fixed for the rest of the study (same backbone, preprocessing,
augmentation and splitting code as every other result reported here
Section~\ref{sec:baseline}). On HandPD spiral with EfficientNet-B0, subject-
grouped evaluation reduces ROC-AUC relative to image-level splitting
(0.991$\to$0.944), and cross-cohort external validation training on one
cohort and testing on the other, with no shared subjects falls further,
to 0.703 (HandPD$\to$Kaggle) and 0.647 (Kaggle$\to$HandPD), both
independently verified under explicit, logged settings
(Section~\ref{sec:limitations}): meaningfully
above chance, but far below the field's headline numbers. The two other
architectures we tested (VGG-16, ResNet-50) showed this same qualitative
pattern but with far higher across-seed variance in the cross-cohort setting
(VGG-16 Kaggle$\to$HandPD AUC std $\pm0.22$), so we anchor the rest of this
study's causal comparisons on EfficientNet-B0 alone. This pattern honest
evaluation reducing apparent accuracy is consistent with subject-identity
leakage as the cause a network can succeed at test time by recognising the
person rather than the disease, because most datasets place several
drawings from the same participant on both sides of an image-level split.

What is not established by this replication, or by prior work on this
data, is \emph{where}, spatially, the identity signal responsible for the
gap actually lives. This matters practically. If subject identity is carried mainly by the background
paper grain, scanning artefacts, the pre-printed guide template then a
masking or normalisation step applied before training could plausibly
remove it without touching the diagnostically relevant trace. If instead
identity is carried by the trace itself an individual's stable,
health-unrelated motor style, or a device/pen/scanner rendering fingerprint
that happens to imprint on the stroke then no purely spatial masking
step can remove it without also removing the very information a screening
model needs.

The same background/foreground question has practical precedent in
adjacent identity \hyp confounding literature \cite{chaibubneto2019,tampu2022}:
segment-level and acquisition-level artefacts are known to inflate
apparent performance in translational EEG and OCT classification. To our
knowledge, no prior work on handwriting-based PD screening has directly
measured, rather than assumed, whether such a background-confined leak
exists in this data, or whether the correct diagnosis is a subtler,
trace-intrinsic one.

\textbf{Contributions.}
\begin{enumerate}
  \item We independently replicate the image-level vs.\ subject-grouped
        evaluation gap and the cross-cohort external-validation ceiling under
        a single protocol held fixed throughout the study, and use it to
        justify anchoring the causal analysis on one stable architecture.
  \item We measure subject re-identification and cohort decodability under
        three masking modes (full / background-only / trace-only) using a
        frozen, untrained feature extractor, directly localising identity
        and acquisition signal within the image rather than inferring it
        from downstream accuracy alone.
  \item We build and validate a resolution-controlled background/template
        suppression pipeline, confirm by direct measurement that it removes
        the background-confined signal it targets, and report the causal
        effect of that removal on honest-protocol ROC-AUC using paired,
        subject-level statistics.
  \item We show the intervention is mechanistically verified but
        outcome-null background suppression does not narrow the honest
        generalisation gap, because the residual signal is substantially
        trace-intrinsic. We rule out image-resolution mismatch as an
        alternative explanation.
  \item We argue this negative-but-explained result is itself
        actionable it forecloses a family of preprocessing-only fixes and
        redirects effort toward trace-level and multimodal mechanisms.
\end{enumerate}

% ======================================================================
\section{Related work}
\label{sec:related}

\subsection{Deep learning for PD handwriting classification}
Reference datasets for this task originate with Pereira et al., who
collected spiral and meander drawings at Botucatu Medical School (HandPD)
and later added dynamic smart-pen signals in the NewHandPD extension
\cite{pereira2015,pereira2016newhandpd}. HandPD and the Kaggle set used
here are not the only handwriting-based PD resources in this literature.
Drotár et al.\ independently collected the PaHaW database (37 PD patients,
38 controls, eight tasks including an Archimedean spiral) with full
kinematic and pressure signal capture rather than static images alone
\cite{drotar2016pahaw}, and Impedovo and Pirlo survey the broader dynamic
handwriting-analysis literature for Parkinson's and Alzheimer's assessment
from a pattern-recognition perspective \cite{impedovo2019}. This study is
restricted to HandPD and Kaggle specifically because they are the two
static-image datasets with a shared drawing task (spiral) across
different acquisition pipelines, which is what the cross-cohort design
requires. PaHaW's kinematic signals are a natural resource for exactly the
dynamic-signal extension we flag in Section~\ref{sec:future}, but are not
directly comparable to the static-image probe applied here. Subsequent
work moved from signal-derived images to ImageNet-pretrained
transfer-learning pipelines fine-tuned on the static drawings. Kamran et
al.\ evaluated six pretrained CNNs (AlexNet, GoogLeNet, VGG-16, VGG-19,
ResNet-50, ResNet-101) across handwriting datasets including HandPD and
report accuracies above 90\% \cite{kamran2021}. Gazda et al.\ used
multiple-fine-tuning transfer learning and reported up to 92.7\% on the
NewHandPD spiral task and 99.22\% on combined data \cite{gazda2022}.
Numerous subsequent papers report 94--100\% on these datasets, often with
heavy augmentation of a few hundred base images, and typically with
evaluation performed at the image level either random $k$-fold over
pooled images or a provided fold structure with subject identity
ignored. To our knowledge, no prior work on handwriting-based PD screening
has directly measured, rather than assumed, whether a resulting
background-confined leak exists in this data, or localised the residual
signal after a proposed mitigation.

\subsection{Identity and acquisition confounding in medical machine learning}
The mechanism this study investigates performance inflation from
identity or acquisition artefacts rather than disease signal is
documented across several adjacent domains, though rarely with the direct,
causal signal-localisation test we apply here. Chaibub Neto et al.\
formalised ``identity confounding'' when repeated measurements per
subject are split at the record level, classifiers exploit person-specific
signatures and performance does not generalise to new people
\cite{chaibubneto2019}. Tampu et al.\ showed that improper (per-image)
splitting of OCT datasets inflates accuracy by 5--30\% relative to
subject-wise splitting \cite{tampu2022}. More broadly, Kapoor and
Narayanan's survey of leakage across ML-based science finds it recurring
in at least 17 fields and 294 papers, and proposes a taxonomy of leakage
types this study's design maps onto directly (our background/template
channel is their ``preprocessing on train and test combined'' type,
localised and tested rather than merely flagged) \cite{kapoor2023}. Saeb
et al.\ make a related evaluation-design point from the opposite
direction cross-validation choices that do not approximate the real
deployment use-case (e.g.\ record-level rather than subject-level folds)
can make an estimate \emph{look} rigorous while still not answering the
clinically relevant question, a caution this study's P1/P2/cross-cohort
protocol separation is designed to take seriously rather than relying on
any single split \cite{saeb2017}. Brookshire et al.\ found that
segment-based holdout strongly overestimates accuracy in translational EEG
relative to subject-based holdout \cite{brookshire2024}. In neuroimaging,
subject leakage inflates connectome-based prediction \cite{rosenblatt2024},
and a broader methodological review of Alzheimer's-disease deep learning
found that studies with confirmed subject-wise splitting reported 66--90\%
accuracy while studies at high risk of leakage claimed 95--99\%
\cite{wen2020,alzheimerscoping2025}. Most of this literature diagnoses
that leakage inflates performance and prescribes subject-wise
splitting as the fix; it does not typically ask, once a mitigation is
applied, whether the mitigation demonstrably removed the signal it targets the question this study's paired probe-then-intervene design is built to
answer.

The present study's design and vocabulary sit closer to the
shortcut-learning literature than to subject-splitting leakage
papers alone. Geirhos et al.\ characterise shortcut learning generally a
network solves a task via a decision rule that exploits spurious,
easy-to-learn correlates of the label rather than its intended causal
features \cite{geirhos2020}. DeGrave et al.\ demonstrate this concretely
for medical imaging radiographic COVID-19 detectors were shown, via
saliency and dataset-provenance analysis, to rely on shortcuts (scanner
type, hospital-specific image properties) rather than pathology
\cite{degrave2021}. Cross-site generalisation failures with the same root
cause are documented directly for pneumonia detection from chest
radiographs \cite{zech2018} and reviewed for COVID-19 imaging models more
broadly \cite{roberts2021}. Badgeley et al.\ show a hip-fracture detector's
performance is driven substantially by confounding hospital- and
scanner-level variables rather than the fracture itself
\cite{badgeley2019}. Kaufman et al.\ give the general data-mining taxonomy
of leakage that this literature instantiates \cite{kaufman2012}, and
Varoquaux and Cheplygina catalogue methodological failure modes
including leakage recurring across medical-imaging machine learning
\cite{varoquaux2022}. Finally, our subject re-identification probe is
methodologically close to work treating re-identifiability itself as a
measurable property of medical images Packh\"auser et al.\ show chest
radiographs carry recoverable patient identity, a privacy-relevant finding
that is the same underlying phenomenon our probe exploits for a different
purpose \cite{packhauser2022}. This study's contribution relative to all
of the above is not the general diagnosis that evaluation design can
manufacture apparent accuracy is now well established but the specific,
causal test of where the responsible signal is spatially located
within the image, and whether a proposed remedy demonstrably removes it.

% ======================================================================
\section{Materials and methods}

\subsection{Datasets}
We use two cohorts common in this literature: HandPD \cite{pereira2015}
(92 subjects, 74 PD / 18 healthy, 368 spiral and 368 meander images,
subject identity recoverable from filenames) and the Kaggle
``Parkinson's Drawings'' set \cite{zham2017} (28 subjects, 15 PD / 13
healthy, 102 spiral and 102 wave images, subject identity recoverable from
the encoded volunteer token). No new data was collected for this study.
The two cohorts also differ in class prevalence (HandPD $80.4\%$ PD;
Kaggle $53.6\%$ PD), which is a further, distinct source of cross-cohort
transfer difficulty independent of any leakage or acquisition-signal
mechanism a decision threshold or operating point learned under one
prevalence does not transfer to the other without recalibration, and
threshold-based metrics (sensitivity, specificity, accuracy) are
correspondingly not comparable across cohorts on face value. This is why
ROC-AUC, which is threshold and to first order prevalence-invariant, is
this study's primary endpoint throughout where we use class-weighted
cross-entropy for cross-cohort training (Section~\ref{sec:training}), that
corrects for training-set imbalance but does not by itself correct for a
prevalence mismatch between training and test cohorts, which remains a
limitation (Section~\ref{sec:limitations}).

We evaluate throughout under three protocols P1 (image-level $k$-fold,
$k{=}10$), P2 (subject-grouped $k$-fold, same backbone/preprocessing/
augmentation/splitting code for every result in this study), and
cross-cohort (train on one cohort's spiral task, test on the other's, no
shared subjects) and report subject-level ROC-AUC as the primary
endpoint. The baseline replication under these three protocols is itself
a result of this study (it establishes the evaluation gap the rest of the
paper investigates the mechanism of) and is reported in
Section~\ref{sec:baseline}, not here.

\subsection{Pipeline overview}
Figure~\ref{fig:architecture} summarises the full study as two lanes sharing
a common front end. Every image is first resolution-standardised and
decomposed into an ink mask (Section~\ref{sec:masking}), which yields the
three views (full / background / trace) used throughout. The
\emph{diagnostic} lane (Section~\ref{sec:probe}) passes all three views
through a frozen, untrained feature extractor to measure subject
re-identification and cohort decodability directly. The causal lane
(Sections~\ref{sec:suppression}--\ref{sec:stats}) instead reconstructs a
suppressed image from the background/trace decomposition, trains a
classifier on the same harness used for raw images, and compares
honest-protocol AUC between the two via a paired, subject-level bootstrap.
The two lanes answer different questions where is the signal
versus does removing it help and Section~\ref{sec:discussion} depends on reading
their outputs together, not on either alone.

Reading the two lanes together yields one of three outcomes, which we
label \textbf{Fork A/B/C} and use throughout the rest of this study:
\begin{description}[leftmargin=1.6em,style=nextline]
\item[Fork A - mechanism confirmed and useful.] The causal lane's paired
      $\Delta$AUC confidence interval lies wholly above zero, \emph{and}
      the diagnostic lane confirms the targeted signal dropped in the
      intervened data. The intervention helped, and we know why.
\item[Fork B - mechanism confirmed, not useful.] The diagnostic lane
      confirms the targeted signal dropped in the intervened data (the
      manipulation worked as intended), but the causal lane's CI includes
      zero. The intervention did what it was supposed to; it just was not
      the (whole) answer.
\item[Fork C - intervention uninterpretable.] The causal lane shows a
      shift in either direction without the diagnostic lane
      confirming the targeted signal actually changed. Any AUC movement
      here cannot be attributed to the intended mechanism and must be
      treated as a possible artefact (e.g.\ a different, unidentified leak,
      or a preprocessing side effect) rather than as support for the
      hypothesis.
\end{description}
This is the specific test that lets us tell a genuine null (Fork B) apart
from a failed manipulation (Fork C): a null causal result is only
interpretable once the diagnostic lane has confirmed the manipulation
happened.

\begin{figure}[H]
\centering
\includegraphics[width=0.84\linewidth]{handwriting_architecture.pdf}
\caption{Pipeline overview. Step 1--3: a shared preprocessing front end
(resolution standardisation, illumination flattening, Sauvola thresholding)
produces the ink mask $M$. Step 4: three masked views of every image, shown
here as real cropped samples from the same HandPD spiral drawing pictured
in full in Figure~\ref{fig:maskpanel} full keeps both the printed
guide and pen stroke, background-only blanks everything the
intensity-only mask $M$ labels as ink (both the guide and the pen stroke
together, since $M$ does not distinguish between them by colour),
\emph{trace-only} keeps only that same combined ink region and blanks
everything else. The
diagnostic lane (left) measures
subject and cohort signal directly with a frozen feature extractor. The
causal lane (right) instead intervenes on the data, retrains the full
classifier, and tests whether the honest-protocol AUC changes. Output chips
show exactly which table or figure each lane feeds.}
\label{fig:architecture}
\end{figure}

\subsection{Positive-control construction}
\label{sec:positive-control-methods}
Before trusting a Fork-B (null) verdict anywhere in this study, we also
need independent evidence that the probe-suppress-retrain pipeline itself
is capable of registering a genuine, known shortcut when one is
deliberately present otherwise a null result could simply mean the
method is blind, not that suppression truly does not help. We construct a
positive control with a ground-truth answer: a synthetic, disease-correlated
marker injected into a mirror copy of HandPD spiral, confined by
construction to background pixels only (a faint corner patch, PD-label
darker than HC-label, never overlapping the ink mask;
\texttt{inject\_marker.py}). Because we built this marker ourselves, we
know in advance what should happen: the diagnostic probe
(Section~\ref{sec:probe}) should detect it as a disease-decoding shortcut
in background-only mode, suppression (Section~\ref{sec:suppression})
should remove it, and this should read as Fork A (or, if the retrain step
does not fully track the diagnostic result, expose a limitation of the
causal lane itself), unlike the Fork-B verdicts we report for the real
acquisition signal. Results are reported in
Section~\ref{sec:positive-control}.

\subsection{Masking modes}
\label{sec:masking}
For each image we define three views: \textbf{full} (the unmodified image),
\textbf{background-only} (the image with the estimated ink/stroke pixels
removed, i.e.\ paper, scanning artefacts and any pre-printed guide), and
\textbf{trace-only} (the estimated ink/stroke pixels alone, i.e.\ the drawn
line). The ink mask is obtained by first flattening uneven illumination via
morphological closing on the grayscale image, then applying Sauvola
adaptive thresholding to the flattened image, followed by connected-component
speck filtering to remove sub-pixel noise.

Let $I$ be the grayscale image and $B=\mathrm{close}_{k}(I)$ its morphological
closing with an elliptical structuring element of size
$k=\max(15,\,\lfloor 0.10\,h\rfloor \mid 1)$, where $\mid$ denotes a
bitwise OR with $1$, used here purely to force the preceding value to the
next odd integer when it is even (OpenCV's morphology and thresholding
kernels require an odd side length. $h$ is the image height). Illumination-flattening divides out the estimated background:
\begin{equation}
I_{\mathrm{flat}}(x,y)=\mathrm{clip}\!\left(\frac{I(x,y)}{\max(B(x,y),1)}\times 255,\ 0,\ 255\right).
\label{eq:flatten}
\end{equation}
The ink mask $M$ is then obtained by Sauvola adaptive thresholding of
$I_{\mathrm{flat}}$ over a local window of side $w=\max(15,\,\lfloor
\min(h,\mathrm{width})/20\rfloor \mid 1)$, with local mean $m(x,y)$ and
standard deviation $\sigma(x,y)$:
\begin{equation}
T(x,y)=m(x,y)\Bigl[1+k_S\Bigl(\tfrac{\sigma(x,y)}{R}-1\Bigr)\Bigr],
\qquad
M=\{(x,y): I_{\mathrm{flat}}(x,y) < T(x,y)\},
\label{eq:sauvola}
\end{equation}
with $k_S=0.2$ and dynamic range $R=128$, followed by removing connected
components smaller than $\max(8,\,2\times10^{-5}|M|)$ pixels. The three
views are then $V_{\mathrm{full}}=I$, $V_{\mathrm{trace}}(x,y)=I(x,y)$ if
$(x,y)\in M$ else white, and $V_{\mathrm{background}}(x,y)=I(x,y)$ if
$(x,y)\notin M$ else white. On HandPD, the pre-printed guide
spiral partially overlaps the pen trace, so the trace-only view there still
retains some template; this split is directional rather than exact, and we
treat within-mode differences smaller than 0.05 as noise rather than signal.

\begin{figure}[H]
\centering
\includegraphics[width=0.85\linewidth]{masking_panel.png}
\caption{Example masking modes for one HandPD spiral (top) and one Kaggle
spiral (bottom). The Sauvola mask $M$ thresholds on intensity alone, not
colour, so on HandPD it captures both the pre-printed guide spiral (dark)
and the patient's pen stroke (blue) together as a single ink region
background-only blanks both together rather than cleanly isolating one
from the other, and trace-only likewise retains both together rather than
isolating the pen. This is the directional-not-exact trace/background
split described above it is a colour confound \emph{within} the
intensity mask, not a case of one colour being selectively retained the
pen-colour-separated correction in Section~\ref{sec:handpd-clean} resolves
it directly. On Kaggle, which has no printed guide, the separation is
clean.}
\label{fig:maskpanel}
\end{figure}

\subsection{Identity and cohort probe}
\label{sec:probe}
For each masking mode we extract 1280-dimensional pooled features from an
EfficientNet-B0 backbone pretrained on ImageNet, with the classification
head removed and no fine-tuning (frozen features the same architecture
family used in the baseline replication (Section~\ref{sec:baseline}) and
causal lane below, but here explicitly untrained on any PD data, so any
signal measured cannot come from disease-specific learning).
We report two quantities. Let $f_i\in\mathbb{R}^{1280}$ be the pooled
feature of drawing $i$, $N$ the total number of drawings in the given mode,
and $\mathrm{subj}(i)$ its subject label.
\begin{itemize}
  \item \textbf{Subject re-identification@1}: for each drawing $i$, its
        nearest neighbour by cosine similarity, excluding itself,
        \begin{equation}
        \hat\jmath(i)=\operatorname*{arg\,max}_{j\neq i}\ \frac{f_i\cdot f_j}{\lVert f_i\rVert\,\lVert f_j\rVert},
        \qquad
        \mathrm{ReID@1}=\frac{1}{N}\sum_{i=1}^{N}\mathbb{1}\bigl[\mathrm{subj}(\hat\jmath(i))=\mathrm{subj}(i)\bigr].
        \label{eq:reid}
        \end{equation}
        Chance is the mean, over $S$ subjects, of $(n_k-1)/(N-1)$ where
        $n_k$ is subject $k$'s drawing count.
  \item \textbf{Cohort decodability}: balanced accuracy of a logistic
        regression classifying HandPD vs.\ Kaggle,
        $\mathrm{BalAcc}=\tfrac{1}{2}\bigl(\mathrm{TPR}+\mathrm{TNR}\bigr)$,
        evaluated under
        group $k$-fold cross-validation grouped by subject (so no subject's
        drawings span both folds), against a chance level of 0.50.
\end{itemize}
Every logistic-regression probe in this study (cohort decodability above
disease and subject-identity probes in
Sections~\ref{sec:background-disease} and~\ref{sec:handpd-clean}) uses
scikit-learn's default solver (\texttt{lbfgs}) and $\ell_2$ penalty with
inverse regularisation strength $C=1$ (scikit-learn's default), raised to
\texttt{max\_iter=2000} to ensure convergence on the 1280-dimensional
frozen features we did not tune $C$ or the penalty type, since the probe's
purpose is to detect the presence of decodable signal at a fixed, declared
setting, not to optimise classification accuracy.
Because frozen ImageNet features are not tuned to handwriting, they tend to
\emph{understate} trace-intrinsic identity signal trace-mode
re-identification reported here is therefore a floor, not a ceiling, on
what a trained model could exploit.

\subsection{Background/template suppression (``2a'')}
\label{sec:suppression}
We label this intervention ``2a'' for consistency with this study's
internal code and released data file names
(\texttt{mech2\_suppression.py}, \texttt{Datasets\_2a/}): it is the first
(``a'') of two variants of our second (``2'') de-leaking mechanism, the
other being an unused, intensity-based alternative
(\texttt{suppress\_template\_intensity()}) not reported in this paper. We
build a suppressed mirror of both cohorts using the same ink-mask
pipeline described above illumination flattening, Sauvola thresholding,
and reconstruction of a clean image retaining only the segmented stroke on
a uniform background. Both cohorts are standardised to the same target
height $h_{\mathrm{target}}=512$\,px before suppression, by scale factor
\begin{equation}
s=\frac{h_{\mathrm{target}}}{h},\qquad (w',h')=\bigl(\mathrm{round}(w\cdot s),\ h_{\mathrm{target}}\bigr),
\label{eq:resize}
\end{equation}
so that differing native resolution between cohorts cannot itself become a
spurious acquisition marker. We manually inspected a sample of suppressed
images from both cohorts to confirm strokes were not thinned or broken and,
where relevant, that the printed guide was handled as intended before
treating any downstream number as meaningful.

\subsection{Training configuration}
\label{sec:training}
Every trained-classifier result in this study the baseline replication
(Section~\ref{sec:baseline}), the causal-lane retraining
(Section~\ref{sec:stats}), and the worked-example checkpoint
(Section~\ref{sec:worked}) uses the same configuration: ImageNet-pretrained
backbone with the classifier head replaced by a 2-way linear layer,
fine-tuned end to end with AdamW (learning rate $10^{-4}$, weight decay
$10^{-4}$), batch size 16, up to 40 epochs with early stopping on
validation AUC (patience 8 epochs, no improvement), input images resized
to $224\times224$ and normalised with ImageNet channel statistics. The
validation set used for this early-stopping criterion differs by
protocol, and we disclose this precisely rather than leaving it
ambiguous. In cross-cohort training, the validation slice is a
subject-grouped $\sim$15\% split carved out of the training cohort
only (never touching the target cohort's data), so checkpoint selection
cannot see anything from the external-validation test set. In the
within-cohort P1/P2/LOSO protocols, by contrast, each fold's held-out test
split is also the validation set used for early stopping on that fold
there is no further inner split of the training data. This means the
reported checkpoint for each within-cohort fold is the epoch with the best
AUC on that fold's own test data, which is a mild form of model
selection on the test set and can inflate the within-cohort AUC estimates
modestly (though it cannot inflate them by picking which data is trained
on, only by picking which epoch's weights are kept). We flag this
explicitly in Section~\ref{sec:limitations} as a methodological caveat on
every within-cohort number (P1, P2, LOSO) in this study, while noting it
does not apply to the cross-cohort comparisons that this study's main
causal claims rest on most heavily. Training-fold
images are augmented with random rotation ($\pm12^{\circ}$), random affine
translation ($\pm5\%$) and scaling ($0.9$--$1.1\times$), and colour jitter
(brightness/contrast $\pm0.1$) augmentation is applied inside the training
fold only, after splitting, so it cannot itself become a leakage channel.
Class-weighted cross-entropy is used in cross-cohort training to counter
HandPD's class imbalance ($\sim$80\% PD). All experiments ran on a single
NVIDIA RTX 3050 (4\,GB) GPU, PyTorch 2.5.1 (CUDA 12.1), scikit-learn,
OpenCV and scikit-image; wall-clock time per (architecture, seed) fold was
under 15 minutes on this hardware.

\subsection{Statistical comparison}
\label{sec:stats}
For the causal test, we used the same evaluation harness as the baseline
replication (Section~\ref{sec:baseline}) unchanged in optimiser,
augmentation and stopping criteria on raw versus suppressed images,
under (i) subject-grouped within-cohort evaluation and
(ii) cross-cohort external validation (train on one cohort, test on the
other, no shared subjects), each over four independently verified, matched
seeds (999, 1001--1003; Section~\ref{sec:limitations}). \textbf{We restrict
this comparison to EfficientNet-B0}, the backbone our baseline replication
(Table~\ref{tab:baseline}) identifies as the stable one for cross-cohort
evaluation; VGG-16 and ResNet-50 showed high across-seed variance in that
setting (VGG-16 Kaggle$\to$HandPD std $\pm0.224$) and are not informative
for a two-condition paired comparison at this sample size.

We report the paired difference in subject-level ROC-AUC via a subject-level
bootstrap. Predictions are first aggregated to one mean probability per
(seed, subject). For seed $s$ with common subject set $\mathcal{U}_s$, we
draw $B=2000$ resamples with replacement $\mathcal{U}_s^{(b)}$ and compute
\begin{equation}
\Delta_s^{(b)}=\mathrm{AUC}\bigl(y_{\mathcal{U}_s^{(b)}},\hat p^{\,\mathrm{supp}}_{\mathcal{U}_s^{(b)}}\bigr)
              -\mathrm{AUC}\bigl(y_{\mathcal{U}_s^{(b)}},\hat p^{\,\mathrm{raw}}_{\mathcal{U}_s^{(b)}}\bigr).
\label{eq:boot}
\end{equation}
Pooling $\{\Delta_s^{(b)}\}$ over all seeds and resamples gives the mean
$\Delta$AUC, a 95\% confidence interval from the 2.5th/97.5th percentiles,
and $P(\Delta\mathrm{AUC}>0)$ as the pooled sample proportion with
$\Delta_s^{(b)}>0$. When a prediction
file contains multiple architectures, they are never pooled within a
subject/seed cell for this statistic each comparison is architecture-pure.

\subsection{Statistical power}
\label{sec:power}
Because the causal comparison rests on 4 seeds and, in the cross-cohort
case, 28 total subjects, an absence of a detected effect is only
informative alongside what effect size the design could have detected. We
approximate each comparison's standard error from its bootstrap 95\% CI
half-width ($\mathrm{SE}\approx(\mathrm{CI}_{\mathrm{hi}}-\mathrm{CI}_{\mathrm{lo}})/(2\times1.96)$)
and report the minimum detectable effect (MDE) at conventional 80\% power,
two-sided $\alpha=0.05$ ($\mathrm{MDE}\approx2.80\times\mathrm{SE}$). This
is an approximation, not a simulation-based power analysis, and is
reported as such it is intended to calibrate how much weight a null
result in Section~\ref{sec:results-causal} can bear, not as a precise
sample-size justification. This MDE framing serves the same ``evidence of
absence versus absence of evidence'' purpose that formal equivalence
testing (two one-sided tests, TOST) is designed to address more
rigorously \cite{lakens2017}: a TOST procedure would require us to specify
a smallest-effect-size-of-interest a priori and would yield a direct
statistical rejection of ``an effect at least that large exists,'' rather
than the MDE's looser ``here is roughly what we could and could not have
detected.'' We did not pre-specify an equivalence bound and so report the
MDE approximation instead a formally pre-registered TOST bound is a
natural improvement for any follow-up study making the same kind of null
claim.

\subsection{Implementation: algorithmic summary}
\label{sec:implementation}
Sections~\ref{sec:masking}--\ref{sec:stats} describe each procedure in
prose and, where the operation is a formula, as an equation. Algorithms
\ref{alg:mask}--\ref{alg:boot} give the same four procedures as pseudocode,
matching the released implementation line for line, so that the exact
order of operations in particular where augmentation and masking
sit relative to the train/test split, which is the detail on which a
leakage study lives or dies is unambiguous.

\begin{algorithm}[H]
\caption{Masked-view construction (\texttt{get\_trace\_mask},
Section~\ref{sec:masking})}
\label{alg:mask}
\begin{algorithmic}[1]
\Require raw image $I$, target height $h_{\mathrm{target}}$
\Ensure views $V_{\mathrm{full}}, V_{\mathrm{background}}, V_{\mathrm{trace}}$
\State $I \gets \Call{StandardiseResolution}{I, h_{\mathrm{target}}}$ \Comment{Eq.~\ref{eq:resize}}
\State $G \gets \Call{Grayscale}{I}$
\State $k \gets \max(15,\ \lfloor 0.10 \cdot \mathrm{height}(G) \rfloor \mid 1)$
\State $B \gets \Call{MorphologicalClose}{G, \text{ellipse}(k)}$
\State $I_{\mathrm{flat}} \gets \Call{Clip}{G / \max(B,1) \times 255,\ 0,\ 255}$ \Comment{Eq.~\ref{eq:flatten}}
\State $w \gets \max(15,\ \lfloor \min(h,\mathrm{width})/20 \rfloor \mid 1)$
\State $T \gets \Call{SauvolaThreshold}{I_{\mathrm{flat}}, w, k_S{=}0.2, R{=}128}$ \Comment{Eq.~\ref{eq:sauvola}}
\State $M \gets \{(x,y) : I_{\mathrm{flat}}(x,y) < T(x,y)\}$
\State $M \gets \Call{RemoveSmallComponents}{M,\ \min\_area}$ \Comment{speck filter}
\State $V_{\mathrm{full}} \gets I$
\State $V_{\mathrm{trace}} \gets I$ with pixels $\notin M$ set to white
\State $V_{\mathrm{background}} \gets I$ with pixels $\in M$ set to white
\State \Return $V_{\mathrm{full}}, V_{\mathrm{background}}, V_{\mathrm{trace}}, M$
\end{algorithmic}
\end{algorithm}

\begin{algorithm}[H]
\caption{Signal-localisation probe (\texttt{leak\_decomposition\_probe.py},
Section~\ref{sec:probe})}
\label{alg:probe}
\begin{algorithmic}[1]
\Require image set $\mathcal{D}$ with (subject, cohort) labels; mode $\in
\{\mathrm{full}, \mathrm{background}, \mathrm{trace}\}$; frozen backbone
$\phi$
\Ensure $\mathrm{ReID@1}$, cohort $\mathrm{BalAcc}$
\For{$i \in \mathcal{D}$}
  \State $(\cdot,\cdot,V_i,\cdot) \gets \Call{Algorithm}{1}(I_i)$ restricted to \textit{mode}
  \State $f_i \gets \phi(V_i)$ \Comment{1280-d pooled features, no fine-tuning}
\EndFor
\For{$i \in \mathcal{D}$}
  \State $\hat\jmath(i) \gets \arg\max_{j\neq i}\ \cos(f_i, f_j)$ \Comment{Eq.~\ref{eq:reid}}
\EndFor
\State $\mathrm{ReID@1} \gets \frac{1}{N}\sum_i \mathbb{1}[\mathrm{subj}(\hat\jmath(i)){=}\mathrm{subj}(i)]$
\State fit logistic regression $\mathrm{cohort} \sim f$ under
       $\mathrm{GroupKFold}(\mathrm{groups}{=}\mathrm{subject})$
\State $\mathrm{BalAcc} \gets \tfrac{1}{2}(\mathrm{TPR}+\mathrm{TNR})$ on held-out folds
\State \Return $\mathrm{ReID@1}, \mathrm{BalAcc}$
\end{algorithmic}
\end{algorithm}

\begin{algorithm}[H]
\caption{Background/template suppression, ``2a'' (\texttt{mech2\_suppression.py},
Section~\ref{sec:suppression})}
\label{alg:suppress}
\begin{algorithmic}[1]
\Require dataset root $\mathcal{R}$ (directory tree of raw images)
\Ensure mirrored root $\mathcal{R}'$ of suppressed images
\For{each image path $p \in \mathcal{R}$ (tree walk)}
  \State $I \gets \Call{Read}{p}$
  \State $(\cdot,\cdot,\cdot,M) \gets \Call{Algorithm}{1}(I)$
  \State $I' \gets$ white canvas with $I$'s pixels at $M$ \Comment{= $V_{\mathrm{trace}}$}
  \State write $I'$ to the corresponding path in $\mathcal{R}'$
\EndFor
\State \Return $\mathcal{R}'$
\end{algorithmic}
\end{algorithm}

\begin{algorithm}[H]
\caption{Paired subject-level bootstrap $\Delta$AUC
(\texttt{compare\_2a\_vs\_baseline.py}, Section~\ref{sec:stats})}
\label{alg:boot}
\begin{algorithmic}[1]
\Require baseline predictions $P_{\mathrm{raw}}$, treatment predictions
$P_{\mathrm{supp}}$ (same architecture, same seeds), bootstrap count $B{=}2000$
\Ensure mean $\Delta$AUC, 95\% CI, $P(\Delta\mathrm{AUC}{>}0)$
\State aggregate $P_{\mathrm{raw}}, P_{\mathrm{supp}}$ to one mean probability per (seed, subject)
\State $\Delta \gets [\,]$
\For{seed $s$ in common seeds}
  \State $\mathcal{U}_s \gets$ subjects present in both $P_{\mathrm{raw}}$ and $P_{\mathrm{supp}}$ at seed $s$
  \For{$b = 1$ \textbf{to} $B$}
    \State $\mathcal{U}_s^{(b)} \gets$ resample $\mathcal{U}_s$ with replacement
    \State append $\mathrm{AUC}(P_{\mathrm{supp}}, \mathcal{U}_s^{(b)}) - \mathrm{AUC}(P_{\mathrm{raw}}, \mathcal{U}_s^{(b)})$ to $\Delta$ \Comment{Eq.~\ref{eq:boot}}
  \EndFor
\EndFor
\State \Return $\mathrm{mean}(\Delta)$, $[\mathrm{percentile}_{2.5}(\Delta), \mathrm{percentile}_{97.5}(\Delta)]$, $\mathrm{mean}(\Delta{>}0)$
\end{algorithmic}
\end{algorithm}

% ======================================================================
\section{Results}

\subsection{Baseline replication: image-level vs.\ subject-grouped and cross-cohort evaluation}
\label{sec:baseline}
Before investigating the mechanism, we establish the evaluation gap
ourselves, under the exact backbone, preprocessing, augmentation and
splitting code used for every other result in this study (P1: stratified
image-level $k$-fold, $k{=}10$; P2: subject-grouped $k$-fold; cross-cohort:
train on one cohort's spiral task, test on the other's, no shared
subjects; five random seeds originally, subsequently four independently
verified matched seeds for the EfficientNet-B0 cells see the table
caption for exactly which cells were re-verified; subject-level ROC-AUC).
Table~\ref{tab:baseline}
reports the result. On HandPD spiral, subject-grouped evaluation reduces
mean AUC relative to image-level splitting for all three architectures, and
cross-cohort external validation falls further for all three. VGG-16 and
ResNet-50 show markedly higher across-seed variance than EfficientNet-B0 in
the cross-cohort setting most strikingly VGG-16 on Kaggle$\to$HandPD,
std $\pm0.224$, wide enough that individual seeds range from near-chance to
near-ceiling which is why the causal comparison in
Section~\ref{sec:stats} anchors on EfficientNet-B0 alone.

\begin{table}[H]
\centering
\caption{Baseline replication, subject-level ROC-AUC (mean $\pm$ SD).
P1 = image-level $k$-fold ($k{=}10$); P2 = subject-grouped $k$-fold
($k{=}10$, with a defensive fallback that reduces $k$ automatically if
the minority class ever has fewer than 10 subject groups this fallback
does not trigger for any number reported in this paper, since every
reported cohort has at least 14 subjects in its smaller class; it exists
only to fail safely on a smaller dataset); cross-cohort = external validation,
no shared subjects. The HandPD spiral row and both EfficientNet-B0
cross-cohort cells are independently verified under explicit, logged
settings (four matched seeds), completing verification of every cell used
in this study's causal comparisons (Table~\ref{tab:causal}). The VGG-16
and ResNet-50 cross-cohort cells are carried over from this study's
original harness output and were not independently re-verified within
this revision (Section~\ref{sec:limitations}).}
\label{tab:baseline}
\resizebox{\textwidth}{!}{%
\begin{tabular}{llcc}
\toprule
Comparison & Model & AUC (P1 / within-cohort) & AUC (P2 / cross-cohort) \\
\midrule
HandPD spiral, P1 vs.\ P2 & EfficientNet-B0 & $0.991\pm0.006$ & $0.944\pm0.005$ \\
\midrule
 & & HandPD$\to$Kaggle & Kaggle$\to$HandPD \\
\cmidrule(lr){3-3}\cmidrule(lr){4-4}
Cross-cohort & VGG-16 & $0.783\pm0.006$ & $0.666\pm0.224$ \\
Cross-cohort & ResNet-50 & $0.601\pm0.029$ & $0.656\pm0.068$ \\
Cross-cohort & EfficientNet-B0 & $0.703\pm0.049$ (verified) & $0.647\pm0.056$ (verified) \\
\bottomrule
\end{tabular}%
}
\end{table}

\begin{figure}[H]
\centering
\includegraphics[width=0.95\linewidth]{baseline_fig.png}
\caption{Cross-cohort AUC by architecture and direction, from
Table~\ref{tab:baseline} (mean $\pm$ SD). EfficientNet-B0 is the only
architecture independently re-verified in this revision (green); VGG-16
and ResNet-50 remain the original, unverified five-seed estimates (blue).
VGG-16 and ResNet-50's wide error bars in the Kaggle$\to$HandPD direction
illustrate why the causal comparison anchors on EfficientNet-B0 alone.}
\label{fig:baseline}
\end{figure}

The HandPD spiral row was corrected during revision: an initial run
under-specified the fold count for this comparison (3 folds rather than
the intended 10), and we replaced it with four independently verified
seeds under explicit, logged settings ($k{=}10$, 40 epochs;
Section~\ref{sec:training}), which gave a consistent, tightly-reproducible
result ($\mathrm{SD}\le0.006$ across seeds) modestly higher than the
original. The same verification was subsequently extended to both
EfficientNet-B0 cross-cohort cells (four matched seeds, identical
settings). HandPD$\to$Kaggle moved only modestly ($0.685\to0.703$, within
the original cell's own standard deviation). Kaggle$\to$HandPD moved
further in absolute terms ($0.703\to0.647$) but the verified estimate
falls well within the original cell's very wide standard deviation
($\pm0.130$, driven by only 28 Kaggle training subjects) and the verified
run's own SD is markedly tighter ($\pm0.056$ vs.\ $\pm0.130$) consistent
with a genuinely noisy, underpowered original estimate rather than a
directional bias from the fold-count bug. Neither cross-cohort correction
changed this study's qualitative conclusions. We disclose all corrections
explicitly rather than silently revising the numbers, and flag in
Section~\ref{sec:limitations} that the remaining cross-cohort cells
(VGG-16 and ResNet-50, both directions), generated by the same original
harness invocation, have not yet received the same independent
re-verification.

Table~\ref{tab:baseline} in fact separates two distinct phenomena that the
introduction's framing risks conflating, and we distinguish them
explicitly from here on. \textbf{Subject-identity leakage} is the
image-level/subject-grouped gap on HandPD spiral is comparatively small
in this design: $0.991\to0.944$, $\Delta\mathrm{AUC}\approx-0.05$.
\textbf{Cross-cohort acquisition/domain shift} is subject-grouped
within-cohort versus external validation on a wholly different cohort
is several times larger: $0.944\to0.703/0.647$,
$\Delta\mathrm{AUC}\approx-0.27$. Both are consistent with the general
diagnosis that evaluation design, not disease signal, drives a large share
of reported accuracy \cite{tampu2022,geirhos2020,degrave2021}, but they are
different mechanisms with different remedies: subject-grouped splitting
addresses the first; it does nothing by construction for the second. The
background/template suppression intervention this study tests targets the
acquisition/domain channel specifically (Section~\ref{sec:suppression}),
not identity leakage as such, and we frame the rest of this study around
that distinction rather than the fused ``leakage'' language common in this
literature. We also note the field's headline 95--99\% figures are
typically reported as accuracy, not AUC; the two are not directly
comparable; where the datasets used here are imbalanced (HandPD $\sim$80\%
PD), accuracy alone would overstate discrimination relative to AUC, which
is why AUC is this study's primary endpoint throughout.

What this replication does not establish is
where, spatially, the signal responsible for the domain-shift
component lives the question the rest of this study addresses.

\subsection{Validating the pipeline: a positive control}
\label{sec:positive-control}
We validate the probe-suppress-retrain pipeline itself using a synthetic,
disease-correlated marker injected into a mirror copy of HandPD spiral
and confined by construction to background pixels only (construction
methodology in Section~\ref{sec:positive-control-methods}); because we
built this marker ourselves,
we know in advance what should happen (Fork A), unlike the Fork-B
verdicts we report for the real acquisition signal below.

\begin{figure}[H]
\centering
\includegraphics[width=0.95\linewidth]{positive_control_fig.png}
\caption{Positive-control summary. Left: background-mode disease
decodability collapses from near-ceiling to chance once the injected
marker is suppressed (step a$\to$b). Right: honest-protocol (P2)
AUC drops only partially, from $0.994$ to $0.972$, still above the
marker-free baseline of $0.944$ (step c) - the residual discussed in
Limitations point Ten.}
\label{fig:positive-control}
\end{figure}

\textbf{Step (a) - the probe detects the injected marker.} On the
raw-marked mirror, a frozen-feature logistic-regression probe
(Section~\ref{sec:probe}) decodes the disease label from
background-only pixels alone with AUC $=0.991$ (balanced accuracy
$0.936$) - consistent with a background-confined shortcut precisely the
size and location we injected, and confirming the probe methodology can
detect such a shortcut when it exists.

\textbf{Step (b) - suppression removes it.} On the suppressed-marked
mirror (the marker mirror pushed through the same background/template
suppression pipeline as every other suppressed image in this study),
background-only disease-decodability collapses to chance AUC $=0.461$
(balanced accuracy $0.500$), a drop of $0.53$ AUC. This is the diagnostic
signature Fork A/B require before a causal result is interpretable at
all, and it confirms the suppression step itself removes a genuine
background shortcut when one is present the same intervention that,
applied to the real cohorts, collapses real background decodability to
chance (Table~\ref{tab:probe-supp}) rather than merely appearing to.

\textbf{Step (c) - does removal show up in honest-protocol AUC?} We then
trained the full classifier (identical harness, five seeds 1--5,
EfficientNet-B0) on the marked mirror before and after suppression, subject
level. Raw-marked: P1 $=1.000$, P2 (subject-grouped) $=0.994\pm0.005$.
Suppressed-marked: P1 $=0.999\pm0.001$, P2 $=0.972\pm0.016$
(Figure~\ref{fig:positive-control}). The injected
marker inflates honest-protocol AUC by roughly $0.994-0.944=0.050$ above
this study's independently verified, marker-free HandPD spiral baseline
(Table~\ref{tab:baseline}). suppression recovers only part of this,
$0.994\to0.972$, a $\Delta\mathrm{AUC}$ of $-0.022$, leaving a residual
$0.972-0.944\approx0.028$ still unaccounted for even though the
background-level probe (step b) shows the marker itself fully collapsed to
chance. We read this as a genuine, if unwelcome, positive-control finding
in its own right: a manipulation that the diagnostic lane confirms worked
completely (background AUC to chance) can still leave a partial,
unexplained residual in the trained classifier's honest AUC, plausibly
because a fine-tuned CNN and a frozen-feature linear probe are not
perfectly matched in what they are sensitive to, or because the
suppression pipeline's reconstruction introduces some other, more subtle
CNN-detectable artefact distinct from the marker itself. Either way, this
means a clean diagnostic-lane result is necessary but may not always be
fully sufficient to guarantee the causal lane's AUC movement tracks it
one-for-one in magnitude a caveat we carry into our reading of the real,
non-marker Fork-B results elsewhere in this study (Section~\ref{sec:limitations}).
It does not, however, undermine the core positive-control conclusion: a
known shortcut produces a large, detectable diagnostic signature and a
large, detectable AUC inflation (step a, c-raw), and suppression measurably
reduces both (step b, c-suppressed) the qualitative pattern this
pipeline is designed to detect. The real acquisition signal tested
elsewhere in this study shows no comparable pattern in either lane, which
is why we read it as a genuine Fork-B null rather than a failure of the
method to register a real effect.

\subsection{Raw images: background is a perfect cohort marker, but identity is not background-confined}
\label{sec:cohort-marker}

\begin{table}[H]
\centering
\caption{Identity and cohort probe on raw images (940 drawings, 120
subjects; chance re-identification@1 $=0.008$, chance cohort balanced
accuracy $=0.50$).}
\label{tab:probe-raw}
\begin{tabular}{lcc}
\toprule
Mode & Subject re-ID@1 & Cohort balanced accuracy \\
\midrule
Full        & 0.327 & 1.000 \\
Background  & 0.251 & 1.000 \\
Trace       & 0.269 & 1.000 \\
\bottomrule
\end{tabular}
\end{table}

Background alone is sufficient to separate the two cohorts perfectly. This
confirms that an acquisition-level marker independent of any disease
signal is present in the data and available to any model, exactly the
kind of confound that would inflate cross-cohort performance for the wrong
reason. However, subject re-identification does not collapse in
background-only mode relative to trace-only mode (0.251 vs.\ 0.269, against
a full-image value of 0.327): identity is carried by both channels, not
confined to the background. This is the point at which the noise-floor
rule stated in Section~\ref{sec:probe} flags a weak a priori case for
background suppression closing the honest gap on its own. We note this
rule was stated in this manuscript before the causal comparison
(Section~\ref{sec:results-causal}) was run, but it was not formally
registered (e.g.\ on OSF or AsPredicted) before analysis began, and we
therefore do not describe it as pre-registered.

\subsection{Controls: label-shuffle and stratified re-identification}
\label{sec:reid-controls}
Two further checks address the retrieval metric itself. \textbf{Label-shuffle
negative control:} we permute the (subject, cohort) labels 30 times,
holding the nearest-neighbour assignment fixed, and recompute
re-ID@1 under each permutation. The empirical shuffle distribution
(mean $0.010\pm0.005$, 95\% range $0.005$--$0.021$) is close to, and
brackets, the analytical chance formula used throughout ($0.007$),
confirming both that the chance formula is accurate and that the observed
re-ID@1 ($0.327$ pooled, full view) is not an artefact of the retrieval
implementation it sits an order of magnitude above the entire empirical
null distribution. \textbf{Stratified re-identification:} because cohort is
perfectly decodable (Table~\ref{tab:probe-raw}), a pooled nearest-neighbour
search over both cohorts effectively never crosses cohort boundaries in
practice, so the pooled chance level understates the correct baseline for
a within-cohort retrieval task. Table~\ref{tab:reid-strat} reports re-ID
stratified by cohort and by cohort$\times$task, each against its own
matched chance level, together with re-ID@5 and mean average precision
(mAP), which are less brittle than re-ID@1 to a single near-miss neighbour.

\begin{table}[H]
\centering
\caption{Subject re-identification, full-colour view, stratified by
cohort and task, each against its own matched chance level. Regenerated
2026-08-26 with \texttt{leak\_decomposition\_probe.py save features} and
\texttt{stratified\_reid.py} to resolve a discrepancy between an earlier
version of this table and Table~\ref{tab:probe-raw} (the pooled row must
equal Table~\ref{tab:probe-raw}'s ``Full'' row by construction; the values
below now do, exactly).}
\label{tab:reid-strat}
\begin{tabular}{lccccc}
\toprule
Stratum & $n$ & re-ID@1 & re-ID@5 & mAP & chance \\
\midrule
Pooled (as Table~\ref{tab:probe-raw}) & 940 & 0.327 & 0.640 & 0.146 & 0.007 \\
HandPD (within-cohort)  & 736 & 0.312 & 0.592 & 0.129 & 0.010 \\
Kaggle (within-cohort)  & 204 & 0.377 & 0.814 & 0.227 & 0.031 \\
HandPD spiral only      & 368 & 0.380 & 0.668 & 0.341 & 0.008 \\
HandPD meander only     & 368 & 0.245 & 0.516 & 0.232 & 0.008 \\
Kaggle spiral only      & 102 & 0.382 & 0.775 & 0.427 & 0.026 \\
Kaggle wave only        & 102 & 0.373 & 0.853 & 0.488 & 0.026 \\
\bottomrule
\end{tabular}
\end{table}

The correctly matched chance levels are $1.1$--$4.2\times$ higher than the
pooled figure, confirming that pooling across cohorts and tasks understated
the correct baseline for a structurally constrained retrieval. The effect
size nonetheless survives this correction by a wide margin in every
stratum: re-ID@1 remains $12$--$46\times$ its stratum-matched chance
level, and mAP the least brittle of the three metrics ranges from
$0.13$ to $0.49$ against chance levels of $0.008$--$0.031$. Identity is
therefore substantially re-identifiable within every cohort and task
individually, not only in the pooled figure.

\subsection{Ruling out a rendering/colour confound}
\label{sec:rendering-invariance}
Table~\ref{tab:probe-raw}'s trace-only view keeps each ink pixel's
original colour and intensity (Algorithm~\ref{alg:mask}, line 11): HandPD
ink is blue-tinted on bluish paper (Figure~\ref{fig:maskpanel}), Kaggle ink
is neutral grey on white. Perfect cohort decodability on that view is
therefore ambiguous by construction it could reflect an acquisition or
motor-style fingerprint on the stroke, or simply that the two cohorts'
pens are different colours, and colour alone is trivial for an
ImageNet-pretrained feature extractor to separate. We resolve this by
re-running the identity/cohort probe on two further renderings of
increasing invariance: \textbf{gray} (grayscale plus per-image contrast
normalisation, removing colour and global brightness) and
\textbf{uniform} (the ink mask skeletonised and re-dilated to a fixed
5\,px width, removing colour, brightness, and stroke width/pressure only the drawn path's geometry survives). We also report cohort
ROC-AUC of the logistic regression's decision function alongside balanced
accuracy, since a metric saturated at $1.000$ in every condition cannot by
itself demonstrate that a manipulation changed anything.

\begin{table}[H]
\centering
\caption{Rendering-invariance ladder, raw images (940 drawings, 120
subjects). color is nominally the same trace-only view as
Table~\ref{tab:probe-raw} but is computed by a separate script
(\texttt{rendering\_invariance\_probe.py}) that, unlike
\texttt{leak\_decomposition\_probe.py} behind Table~\ref{tab:probe-raw},
applies resolution standardisation before masking. this is why its re-ID@1
($0.263$) differs slightly from Table~\ref{tab:probe-raw}'s trace-only
value ($0.269$) rather than matching it exactly a $0.006$ discrepancy
from a methodological difference between scripts, not a new measurement
of a different quantity, and well within this study's declared
($<0.05$) noise floor (Section~\ref{sec:probe}). gray =
grayscale, contrast-normalised. uniform =
skeletonised, fixed-width stroke (geometry only).}
\label{tab:rendering}
\begin{tabular}{lccc}
\toprule
View & Subject re-ID@1 & Cohort AUC & Cohort BalAcc \\
\midrule
color   & 0.263 & 1.000 & 1.000 \\
gray    & 0.248 & 1.000 & 1.000 \\
uniform & 0.221 & 1.000 & 0.997 \\
\bottomrule
\end{tabular}
\end{table}

The result is decisive, and in the opposite direction from the colour-only
explanation: cohort decodability does not collapse under uniform ROC-AUC remains at $1.000$ and balanced accuracy falls only from $1.000$
to $0.997$, even though uniform contains no colour, no
intensity, and no stroke-width information whatsoever, only the traced
path itself. A rendering/colour confound predicts collapse toward chance
under \textit{uniform}; instead the cohort signal survives essentially
intact. We conclude the acquisition-level marker identified in
Table~\ref{tab:probe-raw} is substantially geometric (e.g.\ characteristic
spiral proportions, digitisation/registration conventions, or stroke-path
jitter patterns specific to each acquisition pipeline), not a pixel-colour
artefact, which is the premise the rest of this study's ``trace-intrinsic''
argument depends on. Subject re-identification is more graded: it falls
from $0.263$ to $0.221$ across the ladder, indicating that some
identity-relevant signal is colour/rendering-related, but a large residual
(30$\times$ chance) survives into the pure-geometry view as well, so
identity too is not primarily a colour artefact. Because cohort AUC itself
saturates at $1.000$ even under the strictest rendering control, a fully
graded comparison across conditions would additionally benefit from a
continuous separability measure such as log-loss or decision-margin
distributions; we did not compute these here and flag it in
Section~\ref{sec:limitations} as a remaining refinement, though the
balanced-accuracy and AUC evidence already available is sufficient to
reject the colour-confound explanation.

\subsection{Suppression works as intended, but the residual signal is trace-intrinsic}

\begin{table}[H]
\centering
\caption{Identity and cohort probe after background/template suppression
(940 drawings, 120 subjects matching Table~\ref{tab:probe-raw} exactly;
see Section~\ref{sec:limitations} on an earlier file-count discrepancy that
this corrected run resolves).}
\label{tab:probe-supp}
\begin{tabular}{lcc}
\toprule
Mode & Subject re-ID@1 & Cohort balanced accuracy \\
\midrule
Full        & 0.236 & 0.974 \\
Background  & 0.005 & 0.500 \\
Trace       & 0.236 & 0.974 \\
\bottomrule
\end{tabular}
\end{table}

Suppression collapses both background-mode subject re-identification and
background-mode cohort decodability to chance (0.005 and 0.500
respectively), confirming the intervention removed exactly the signal it
targeted. Full-image and trace-only values are now identical (0.236),
meaning all remaining identity signal, after suppression, is trace-borne.
Cohort decodability through the trace channel remains high (0.974), and
Section~\ref{sec:rendering-invariance2} shows this is not a rendering
artefact: re-running the rendering-invariance ladder on these suppressed
images gives cohort ROC-AUC $1.000\to1.000\to1.000$ across
color/gray/uniform renderings an acquisition fingerprint, consistent
with device, pen, or scan-rendering differences between cohorts, survives
suppression because it is imprinted on the stroke's geometry itself, not
only on the surrounding paper or template, and not on pixel colour.

\subsection{Confirming the geometric result on suppressed images}
\label{sec:rendering-invariance2}
We repeated the rendering-invariance ladder (Section~\ref{sec:rendering-invariance})
on the suppressed images themselves, to check that the trace-channel
cohort signal in Table~\ref{tab:probe-supp} is not itself a residual
rendering artefact of the suppression pipeline. We note, correcting an
earlier draft of this manuscript, that suppression does \emph{not} retain
the original ink colour at masked pixels: \texttt{mech2\_suppression.py}
renders every retained ink pixel as pure black on white regardless of its
source colour, which is also why the \textit{color} and \textit{gray}
rows of Table~\ref{tab:rendering-supp} below are numerically identical --
there is no colour information left in these images for the two renderings
to differ on. This is itself a useful internal check rather than a defect:
any residual cohort or identity signal recovered from the ``color'' row of
suppressed images is, by construction, purely a geometry/intensity effect,
not a true colour effect.

\begin{table}[H]
\centering
\caption{Rendering-invariance ladder, suppressed images (940 drawings, 120
subjects).}
\label{tab:rendering-supp}
\begin{tabular}{lccc}
\toprule
View & Subject re-ID@1 & Cohort AUC & Cohort BalAcc \\
\midrule
color   & 0.236 & 1.000 & 0.999 \\
gray    & 0.236 & 1.000 & 0.999 \\
uniform & 0.221 & 1.000 & 0.997 \\
\bottomrule
\end{tabular}
\end{table}

\begin{figure}[H]
\centering
\includegraphics[width=0.95\linewidth]{rendering_ladder_fig.png}
\caption{Rendering-invariance ladder, raw (Table~\ref{tab:rendering}) vs.\
suppressed (Table~\ref{tab:rendering-supp}) images. Cohort ROC-AUC (blue)
stays pinned at $1.000$ across color$\to$gray$\to$uniform in both panels
the signature a colour/rendering confound would not produce while
subject re-ID@1 (orange) declines only gradually, confirming a large
geometric residual survives even the strictest rendering control.}
\label{fig:rendering-ladder}
\end{figure}

The pattern is the same as on raw images (Table~\ref{tab:rendering}) and,
if anything, cleaner: cohort ROC-AUC does not move at all across the
ladder, and balanced accuracy falls only from $0.999$ to $0.997$ under the
fully geometry-only rendering. Combined with Table~\ref{tab:probe-supp},
this rules out both a background/template confound (already shown by
suppression itself) and a colour/rendering confound (shown here) as
explanations for the cross-cohort signal that survives in the trace. What
remains is the drawn path's geometry.

\subsection{Suppression does not move the honest-protocol AUC}
\label{sec:results-causal}

\begin{table}[H]
\centering
\caption{Paired subject-level bootstrap comparison, raw vs.\ suppressed
images, EfficientNet-B0. All three rows are independently verified under
explicit, logged settings (four matched seeds: 999, 1001--1003; $k{=}10$,
40 epochs, Section~\ref{sec:training}), replacing the earlier five-seed
(seeds 1--5) results from this study's original harness invocation
(Section~\ref{sec:limitations}).}
\label{tab:causal}
\resizebox{\textwidth}{!}{%
\begin{tabular}{lccc}
\toprule
Comparison & Raw AUC & Suppressed AUC & Paired $\Delta$AUC (95\% CI) \\
\midrule
Cross-cohort HandPD$\rightarrow$Kaggle (verified) & 0.703 & 0.662 &
  $-0.040$ ($-0.339$ to $+0.221$) \\
Cross-cohort Kaggle$\rightarrow$HandPD (verified) & 0.647 & 0.664 &
  $+0.017$ ($-0.404$ to $+0.309$) \\
Within-cohort, HandPD spiral (verified) & 0.944 & 0.898 &
  $-0.045$ ($-0.134$ to $+0.029$) \\
\bottomrule
\end{tabular}%
}
\end{table}

All three confidence intervals include zero ($P(\Delta\text{AUC}>0)=0.396$,
$0.723$, and $0.121$ respectively), so none reaches the pre-specified
threshold for a supported improvement. Applying the approximate
minimum-detectable-effect calculation of Section~\ref{sec:power} to these
now-verified confidence intervals, this
design could reliably detect ($80\%$ power, two-sided
$\alpha{=}0.05$) an effect of roughly $\mathrm{MDE}\approx0.40$ AUC
(HandPD$\to$Kaggle), $0.51$ AUC (Kaggle$\to$HandPD), and $0.12$ AUC
(within-cohort) both cross-cohort MDEs widened relative to the original,
under-specified run's own CIs (previously $\approx0.36$ and $\approx0.37$),
confirming both directions are less powered than originally reported, not
more. \textbf{A null result at this power supports only ``we did
not detect an improvement of this size or larger''}; it does not support
``suppression does not help'' as a general claim, and we revise our
language accordingly for the rest of this section and the abstract. The
three comparisons are not
equally informative, however. Both cross-cohort directions give wide
intervals roughly symmetric around a near-zero point estimate ($-0.040$ and
$+0.017$) consistent with a genuinely underpowered null in either
direction, given only 28 Kaggle subjects total. Both cross-cohort rows are
now independently verified on four matched seeds. each moved only modestly
from its original five-seed estimate ($-0.009\to-0.040$ for
HandPD$\to$Kaggle, $+0.022\to+0.017$ for Kaggle$\to$HandPD), well within
the width of either interval, and the Kaggle$\to$HandPD interval widened
further under verification ($\pm0.130\to\pm0.191$ on the suppressed side
alone), underscoring how underpowered this direction is at only 28 test
or training subjects. The within-cohort row, now
independently verified on four matched seeds under explicit, logged
settings, gives a narrower interval centred modestly negative but clearly
spanning zero ($P(\Delta\text{AUC}>0)=0.121$, a one-sided bootstrap resample
proportion, not a posterior probability no prior is specified); this
softens the original five-seed result's borderline signal
($P(\Delta\text{AUC}>0)=0.054$) into an unambiguous Fork-B null. We do not
treat this as evidence that suppression cost genuine signal: the interval
includes zero, and the point estimate itself is small. We report it as a
supported null rather than an unresolved possibility now that the matched
raw/suppressed pair for these seeds is independently verified.
Section~\ref{sec:power} quantifies the minimum effect this design could
have detected.

\begin{figure}[H]
\centering
\includegraphics[width=0.75\linewidth]{auc_forest.png}
\caption{Paired $\Delta$AUC (suppressed $-$ raw), EfficientNet-B0 only,
subject-level bootstrap 95\% CI. All three rows use four independently
verified matched seeds (999, 1001--1003). All three intervals include
zero; both cross-cohort directions are wide and centred near zero
(underpowered null in either direction, given only 28 Kaggle subjects
total), the within-cohort interval is narrower and modestly negative but
clearly spans zero ($P(\Delta\text{AUC}>0)=0.121$).}
\label{fig:forest}
\end{figure}

\subsection{Ruling out resolution and aspect-ratio artefacts}
\label{sec:resolution-artefact}
Because unequal native resolution between cohorts can itself masquerade as
a cohort marker, we verified directly that both cohorts were held at the
same 512\,px target height throughout suppression (HandPD images
$\approx$520--555$\times$512\,px, aspect ratio $\approx0.99$--$1.08$.
Kaggle images exactly 512$\times$512\,px, aspect ratio exactly $1.0$). This
residual aspect-ratio difference is itself a further, cohort-correlated
geometric variable worth controlling every image is subsequently resized
to a forced $224\times224$ square for the classifier and probe backbone
(Section~\ref{sec:probe}), which stretches HandPD images by a small,
cohort-specific horizontal factor that Kaggle images never experience,
since Kaggle is already exactly square. To test whether this stretch, not
genuine stroke geometry, explains the geometry-only cohort signal
(Table~\ref{tab:rendering}), we re-ran the \textit{uniform} condition with
images padded to a square canvas (preserving aspect ratio) before the
$224\times224$ resize, instead of stretched directly. Cohort AUC was
unchanged ($1.000$ under both stretch and pad-to-square) and balanced
accuracy was, if anything, marginally higher under the aspect-preserving
control ($0.9975$ vs.\ $0.9969$) the opposite of what a stretch
artefact would predict. Neither the resolution-matching step nor the
downstream forced-square resize explains the trace-channel cohort signal
reported above.

\subsection{A genuinely clean pen-only split on HandPD, and a corrected colour assignment}
\label{sec:handpd-clean}
Section~\ref{sec:masking} already discloses that HandPD's trace-only view
is not a clean separation the intensity-only Sauvola mask
(Algorithm~\ref{alg:mask}) keeps both the printed guide and the pen stroke
as ``ink,'' and the two are different colours. An external review of this
manuscript identified that our colour assignment for which colour
is which was backwards throughout this study's internal code and text
(including an earlier draft of this section, which described the guide as
blue and the pen as dark/neutral). We verified the correct assignment
directly rather than taking either the original assumption or the
reviewer's inference on faith overlaying the high and low-blue-dominance
masks across eight different HandPD subjects' raw images (after
resolution standardisation) shows the low-blue-dominance (dark) mask
overlays into a crisp, near-identical spiral across different people's
sheets (cross-subject mean IoU $0.225$), exactly what a single printed
template produces, while the high-blue-dominance (blue) mask overlays into
a thick, diffuse band (IoU $0.104$), consistent with hand-drawn variation.
Direct visual inspection of raw images confirms this unambiguously: the
\textbf{black/dark spiral is the printed guide} (thin, smooth,
geometrically regular) and the \textbf{blue spiral is the patient's actual
pen stroke} (thick, jagged, tremulous). We correct this throughout the
manuscript and the released code (\texttt{handpd\_clean\_split.py}), and
disclose the error explicitly rather than silently revising affected
numbers.

Using the corrected rule (keep high-blue-dominance pixels within the ink
mask the pen; drop low-blue-dominance pixels the guide), $47.0\%$ of
the pixels the intensity-only mask had labelled ``ink'' are guide, not pen
($11.0\%$ of the image on average, down to $5.8\%$ pen-only) essentially
the same proportion reported before the colour-direction fix, since that
number depends only on the blue-dominance threshold splitting the mask
into two groups, not on which group we call ``pen.'' What changes
substantively is which number is the real pen-only re-identification
result. Re-running subject re-identification on HandPD alone under the
\emph{corrected} pen-only view gives re-ID@1 $=0.236$ (chance $0.0095$),
compared with $0.261$ on the original guide-contaminated trace (both
measured on the same 736 HandPD drawings). The earlier, backwards
computation (re-ID@1 $=0.072$) was in fact isolating the \emph{guide},
and its near-chance-adjacent value is exactly what we would expect from a
machine-printed template that carries little genuine subject signal
a useful sanity check on the correction, not a contradiction of it. With
the correction applied, HandPD's trace-intrinsic identity signal is
not substantially reduced by removing the guide ($0.261\to0.236$,
retaining $90\%$ of the original re-ID rate) the signal was predominantly
genuine, pen-driven, and the previous ``partially confirmed, overstated''
characterisation is itself corrected the trace-intrinsic claim is
better supported than the original (uncorrected) analysis suggested, not
worse.

This correction also directly answers the sharper concern it raises: if
HandPD's ink mask contains two overlaid spirals (guide and pen) while
Kaggle's contains only one (Kaggle has no printed guide at all), could the
perfect cohort decodability under geometry-only rendering
(Table~\ref{tab:rendering}) be trivially explained by ``two spirals
present vs.\ one,'' with nothing to do with any genuine acquisition
fingerprint in the stroke? We re-ran the full rendering-invariance ladder
(Section~\ref{sec:rendering-invariance}) with HandPD's guide correctly
excluded (Kaggle unchanged, since it has no guide to remove) and report
the result as Table~\ref{tab:rendering-penonly}.

\begin{table}[H]
\centering
\caption{Rendering-invariance ladder, HandPD guide correctly excluded
(raw images, 92 HandPD subjects pen-only + all Kaggle subjects, 2026-10
re-run following external review). Cohort AUC remains at $1.000$ across
every rendering, directly refuting the two-spirals-vs-one-spiral
explanation for cohort decodability.}
\label{tab:rendering-penonly}
\begin{tabular}{lccc}
\toprule
View & Subject re-ID@1 & Cohort AUC & Cohort BalAcc \\
\midrule
color   & 0.244 & 1.000 & 1.000 \\
gray    & 0.251 & 1.000 & 0.990 \\
uniform & 0.202 & 1.000 & 0.981 \\
\bottomrule
\end{tabular}
\end{table}

Cohort ROC-AUC stays perfectly at $1.000$ under every rendering, including
\textit{uniform} (pure stroke geometry, guide excluded) the same
decisive pattern as the original, guide-contaminated ladder
(Table~\ref{tab:rendering}). If the two-spirals confound explained the
original result, removing one of the two spirals should have moved cohort
AUC toward chance; it did not move at all. We repeated this on the
suppressed images (constructing the ``suppressed, pen-only'' view directly
from the raw pen-only mask via the same black-on-white rendering
convention as Section~\ref{sec:suppression}'s pipeline, because we also
discovered that the existing suppressed images retain no colour
information at all see the correction below) color/gray/uniform give
cohort AUC $1.000$/$1.000$/$1.000$ and re-ID@1 $0.239$/$0.239$/$0.202$,
again matching the original pattern closely. We conclude the guide-overlap
confound this correction set out to test does not explain this study's
central cohort-decodability finding; the signal survives being properly
isolated to the pen alone.

\textbf{A second, independent correction.} While constructing the
suppressed-image pen-only view above, we found that the suppressed images
in \texttt{Datasets\_2a/} contain only pure black (0,0,0) and pure white
(255,255,255) pixels no intermediate colour or intensity survives
suppression at all. This contradicts an earlier claim elsewhere in this
manuscript (Section~\ref{sec:rendering-invariance2}) that the suppression
pipeline ``retain[s] the original ink colour at masked pixels.'' It does
not; \texttt{mech2\_suppression.render()} renders every ink pixel as pure
black regardless of its original colour. We correct that claim at its
source (Section~\ref{sec:rendering-invariance2}) and disclose it here
rather than silently revising it.

\subsection{Background-only disease decodability, and whether suppression narrows the identity-leakage gap}
\label{sec:background-disease}
Section~\ref{sec:cohort-marker} establishes that background (template and paper)
is a near-perfect cohort marker. This leaves open a sharper,
within-cohort question that a reviewer of this manuscript raised directly
is the background also predictive of disease status within a single
cohort, where it cannot simply be riding on the already-perfect
cohort signal? If background carries no disease information, the honest
finding that suppression does not move honest-protocol AUC
(Section~\ref{sec:results-causal}) is close to guaranteed by construction
and not very informative. If background does carry disease
information, the same null becomes a stronger result: it shows that
suppressing a genuinely disease-predictive channel still does not move the
honest classifier's accuracy, because the classifier was not relying on
that channel in the first place.

We ran a frozen-feature probe for within-cohort disease decodability
(PD vs.\ healthy control) under \textit{full}, \textit{background-only},
and \textit{trace-only} views, separately for each cohort (pooling cohorts
would reintroduce the cohort confound) and separately for raw and
suppressed images (Table~\ref{tab:background-disease}).

\begin{table}[H]
\centering
\caption{Within-cohort disease (PD vs.\ HC) decodability by masking mode,
raw vs.\ suppressed images (frozen EfficientNet-B0 features, GroupKFold
logistic regression probe).}
\label{tab:background-disease}
\begin{tabular}{llccc}
\toprule
Cohort & Condition & Full AUC & Background AUC & Trace AUC \\
\midrule
HandPD & raw        & 0.926 & 0.921 & 0.893 \\
HandPD & suppressed & 0.894 & 0.461 & 0.894 \\
Kaggle & raw        & 0.891 & 0.923 & 0.875 \\
Kaggle & suppressed & 0.875 & 0.253 & 0.875 \\
\bottomrule
\end{tabular}
\end{table}

On raw images, background-only disease decodability is high in both
cohorts (AUC $0.921$ HandPD, $0.923$ Kaggle) essentially as high as the
full image directly refuting the possibility that background is only a
cohort shortcut with no disease information. Background is, on these raw
images, itself a near-full disease marker, almost certainly via
PD-associated differences in paper handling, scanning, or acquisition
protocol rather than anything about the stroke. Suppression collapses this
to chance or below in both cohorts ($0.461$ HandPD, $0.253$ Kaggle
the Kaggle value is likely small-sample noise given $n=102$, but is, if
anything, evidence the channel is fully removed rather than only
weakened), confirming the intervention worked exactly as intended on a
real and substantial disease-predictive confound, not only on a synthetic
marker (Section~\ref{sec:positive-control}) or on cohort identity
(Section~\ref{sec:cohort-marker}). Trace-only AUC is essentially unchanged by
suppression within each cohort ($0.893\to0.894$ HandPD, $0.875\to0.875$
Kaggle), consistent with trace and background being probed independently.
That the honest-protocol AUC still does not move after suppressing a
background channel this strongly disease-predictive
(Section~\ref{sec:results-causal}) is therefore a materially stronger
null result than if background had carried little disease signal to begin
with the classifier was leaving a genuinely informative channel on the
table, not merely one that happened to carry no signal.

We also address a related question raised about this study's title claim:
does suppression narrow the gap between image-level (P1, no grouping) and
subject-grouped (P2) honest-protocol AUC the identity-leakage gap the
paper is centrally about? Table~\ref{tab:probe-raw}'s companion causal runs
already contain both P1 and P2 rows for the suppressed condition; no new
training was required to extract them. Raw images give P1 $=0.991\pm0.006$
and P2 $=0.944\pm0.005$ (gap $0.047$); suppressed images give P1
$=0.961\pm0.021$ and P2 $=0.898\pm0.019$ (gap $0.063$, four seeds).
Suppression does not narrow the P1-P2 gap; if anything it is nominally
larger, though the two gaps are not statistically distinguishable at this
sample size. This is consistent with the rest of this study's results the
identity-leakage gap and the background/template channel are not the same
mechanism, and removing the latter does not resolve the former.

\subsection{Calibration collapses under cross-cohort shift}
\label{sec:calibration}
AUC measures ranking quality, not whether predicted probabilities are
trustworthy a clinical screening tool needs both. We report the Brier
score and expected calibration error (ECE, 10 equal-width bins) for
EfficientNet-B0's subject-level predicted probabilities under P1, P2, and
both cross-cohort directions (mean $\pm$ SD over five seeds), reusing the
predictions already generated for Table~\ref{tab:baseline}. The P1/P2 rows
therefore carry the same provenance caveat as Table~\ref{tab:baseline}'s
original (pre-correction) predictions (Section~\ref{sec:limitations}) we
report them as an illustrative comparison of within-cohort versus
cross-cohort calibration, not as a precision claim on the exact values.

\begin{table}[H]
\centering
\caption{Calibration, EfficientNet-B0, HandPD spiral (five seeds). Lower
is better for both metrics.}
\label{tab:calibration}
\begin{tabular}{lcc}
\toprule
Regime & Brier score & ECE \\
\midrule
P1 (image-level)        & $0.090\pm0.003$ & $0.141\pm0.012$ \\
P2 (subject-grouped)    & $0.100\pm0.009$ & $0.106\pm0.015$ \\
Cross-cohort HandPD$\to$Kaggle & $0.275\pm0.017$ & $0.239\pm0.014$ \\
Cross-cohort Kaggle$\to$HandPD & $0.257\pm0.029$ & $0.330\pm0.037$ \\
\bottomrule
\end{tabular}
\end{table}

\begin{figure}[H]
\centering
\includegraphics[width=0.85\linewidth]{calibration_fig.png}
\caption{Brier score and ECE from Table~\ref{tab:calibration} (mean $\pm$
SD). Both metrics roughly double to triple moving from within-cohort
(P1/P2) to cross-cohort evaluation, visualising that domain shift degrades
not just discrimination but the trustworthiness of the predicted
probabilities themselves.}
\label{fig:calibration}
\end{figure}

Both metrics are similar under P1 and P2 honest subject-grouped
splitting costs some discrimination (Table~\ref{tab:baseline}) but does
not by itself substantially degrade calibration and then roughly
double to triple under cross-cohort transfer in both directions. This
matches this study's motivation for evaluating cross-cohort transfer at
all (Section~\ref{sec:baseline}) the dominant failure mode under domain shift
is not merely lower discrimination but the predicted probabilities
becoming actively unreliable, which matters directly for any clinical
operating point set on one cohort and deployed on another. We do not
recalibrate (e.g.\ Platt scaling or temperature scaling on a held-out
cross-cohort sample) here doing so, and reporting sensitivity/specificity
at a clinically motivated operating point per cohort, is a natural
extension we leave to future work.

\subsection{Generalisation across drawing tasks}
\label{sec:othertasks}
The causal comparison above uses the spiral task because it is the only
task name shared by both cohorts (HandPD provides spiral and meander,
Kaggle provides spiral and wave), so it is the only task on which
cross-cohort external validation is structurally possible. Within-cohort
comparisons are not subject to that constraint, and we repeated the
identical baseline-replication and suppression pipeline (same backbone,
same training configuration, Section~\ref{sec:training}) on HandPD's
other task (meander) and Kaggle's other task (wave). Table~\ref{tab:othertasks}
shows the same qualitative pattern holds: image-level splitting inflates
apparent AUC over subject-grouped splitting on both tasks, and background
suppression again produces a paired $\Delta$AUC whose CI includes zero.
This is not a new mechanism finding it is the spiral-task result
replicated on two further tasks but it addresses whether Table~\ref{tab:causal}'s
conclusion is specific to spiral drawings, and the answer is no. As with
Table~\ref{tab:calibration}, these five-seed numbers carry the same
provenance caveat as Table~\ref{tab:baseline}'s original, pre-correction
predictions (Section~\ref{sec:limitations}) they were generated by the
original harness invocation, not independently re-verified under the
explicit, logged four-seed protocol used for the HandPD spiral task and
both cross-cohort directions. We report them as an illustrative
replication of the qualitative pattern on two further tasks, not as a
precision claim on the exact values, and flag independent re-verification
of this table as a natural, lower-priority follow-up alongside the
VGG-16/ResNet-50 cross-cohort cells (Section~\ref{sec:limitations}).

\begin{table}[H]
\centering
\caption{Within-cohort generalisation to HandPD meander and Kaggle wave,
EfficientNet-B0 (five seeds).}
\label{tab:othertasks}
\resizebox{\textwidth}{!}{%
\begin{tabular}{lccc}
\toprule
Task & AUC (P1) & AUC (P2) & Suppression $\Delta$AUC (95\% CI) \\
\midrule
HandPD meander & $0.929\pm0.013$ & $0.911\pm0.022$ & $-0.053$ ($-0.167$ to $+0.051$) \\
Kaggle wave     & $0.753\pm0.022$ & $0.627\pm0.056$ & $+0.013$ ($-0.211$ to $+0.267$) \\
\bottomrule
\end{tabular}%
}
\end{table}

Note also that the frozen-feature identity/cohort probe
(Section~\ref{sec:probe}, Tables~\ref{tab:probe-raw}-\ref{tab:probe-supp})
already pools every available task per cohort by construction ink-mask
extraction (Algorithm~\ref{alg:mask}) does not depend on which task a
drawing belongs to so those two tables already reflect spiral and
meander together for HandPD (736 drawings $=$ 368+368) and spiral and wave
together for Kaggle (204 drawings $=$ 102+102), not spiral alone.

\subsection{Worked example: one prediction under each masking mode}
\label{sec:worked}
Tables~\ref{tab:probe-raw}-\ref{tab:causal} report aggregate statistics
to make the underlying phenomenon concrete, we trained one full classifier
(EfficientNet-B0, HandPD spiral, subject-grouped split released as a
runnable demonstration, see Data and code availability) and applied it to
a single held-out drawing from a PD subject under all three masking
modes, using exactly the preprocessing described in
Section~\ref{sec:masking}. \textbf{This checkpoint is a single fold from a
single split (fixed seed), not the multi-seed mean reported in
Table~\ref{tab:baseline}}, and its held-out image-level AUC of $0.996$ on
that one fold's 76-image validation set is accordingly not directly
comparable to though, following this study's own baseline correction
(Section~\ref{sec:baseline}), now much closer to Table~\ref{tab:baseline}'s
$0.944\pm0.005$ four-seed mean it is shown to illustrate the masking-mode
behaviour of one representative trained model, not as an additional
performance claim.
Figure~\ref{fig:worked} shows the
result under the full image and the trace-only view, the model predicts
PD with high confidence ($P(\mathrm{PD})=0.996$ and $0.988$
respectively) under the background-only view with the drawn stroke
itself removed the prediction collapses to near chance
($P(\mathrm{PD})=0.487$) and flips to the wrong class. This is a single
illustrative case, not a statistical claim, but it makes Table~\ref{tab:probe-raw}'s
aggregate finding tangible this model's PD-relevant evidence for this
drawing was not carried by the background.

\begin{figure}[H]
\centering
\includegraphics[width=0.95\linewidth]{worked_example.png}
\caption{One HandPD spiral from a PD subject, classified under each
masking mode by the trained demonstration checkpoint (a single fold, held-out
image-level AUC $0.996$ on that fold not the Table~\ref{tab:baseline}
four-seed mean). The prediction survives background removal (trace-only,
$P(\mathrm{PD})=0.988$) but collapses toward chance and flips to healthy
control when only the background is shown ($P(\mathrm{PD})=0.487$),
concretely illustrating the aggregate pattern in
Table~\ref{tab:probe-raw}-\ref{tab:probe-supp}.}
\label{fig:worked}
\end{figure}

% ======================================================================
\section{Discussion}
\label{sec:discussion}

Taken together, these results answer the localisation question our own
baseline replication (Section~\ref{sec:baseline}) leaves open, and the
answer is not the convenient one. The
background does carry a strong, disease-irrelevant acquisition marker
(Table~\ref{tab:probe-raw}), which is exactly the kind of leak a
preprocessing step ought to be able to remove and our suppression
pipeline demonstrably does remove it (Table~\ref{tab:probe-supp}). But
removing it does not recover the honest generalisation gap
(Table~\ref{tab:causal}), because the same acquisition information, along
with subject identity, is also encoded in the trace itself. A model does
not need the paper or the template to tell HandPD and Kaggle drawings
apart, or to tell two different people's drawings apart it can do both
from the stroke alone. If anything, the within-cohort comparison leans
mildly negative rather than perfectly null (Table~\ref{tab:causal}),
suggesting suppression can cost a small amount of genuine signal alongside
the leak it removes a
caution against treating background masking as a free intervention even
when it is not the answer to the generalisation gap.

This has a direct methodological implication for the field masking-based
or background-normalisation-based de-leaking strategies should not be
assumed to work without the kind of paired, mechanism-verified test applied
here. A positive AUC shift after such an intervention is not sufficient
evidence that leakage was removed, and as importantly a null AUC
shift, on its own, is not sufficient evidence that the intervention failed
to work as a manipulation. Distinguishing those two failure modes requires
the same signal-localisation probe used here, applied both before and
after the intervention, not only a before/after accuracy comparison.

We deliberately report this as a negative result rather than omitting it
or reframing the study post hoc. The literature on data leakage in medical
imaging has repeatedly found that a large share of headline accuracy is an
artefact of evaluation design rather than of learned disease signal
\cite{tampu2022} this study adds a specific, falsifiable finding to that
picture for handwriting-based PD screening, at least one plausible and
mechanistically verified remediation background/template suppression
is not the answer, and the reason is now localised rather than speculative.

\subsection{Implications for future mechanisms}
\label{sec:future}
Because the residual signal is trace-intrinsic, the productive next step is
not further background engineering but methods that either (i) make the
trained representation invariant to the acquisition-specific component of
the trace while preserving the disease-relevant component (e.g.\
domain-adversarial \cite{ganin2015} or correlation-alignment
\cite{sun2016coral} training between cohorts), or
(ii) supply information the static trace image does not fully capture,
such as raw dynamic pen signal (pressure, velocity, in-air time), which may
separate motor-disease signal from rendering-fingerprint signal more
cleanly than a static image can. We treat (ii) as motivated future work
rather than a result of the present study, but we did run a small pilot of
one instance of
(i) Deep CORAL \cite{sun2016coral}, which adds a differentiable
penalty aligning source/target feature covariance during training, reusing
this study's exact backbone and cross-cohort harness so the comparison is
apples-to-apples with Table~\ref{tab:baseline}'s plain baseline and
Table~\ref{tab:causal}'s suppression result (HandPD$\to$Kaggle, same four
verified seeds, EfficientNet-B0). CORAL AUC $=0.622\pm0.049$, \emph{below}
both the plain baseline ($0.703\pm0.049$) and the suppressed-image result
($0.662\pm0.074$) for this direction (Figure~\ref{fig:interventions}),
though all three estimates' wide,
overlapping uncertainty (Section~\ref{sec:power}) means we do not read this
as CORAL actively hurting generalisation, only as failing to demonstrate
the improvement the domain-invariance hypothesis would predict, on this
one direction, at this one hyperparameter setting (default CORAL weight,
no sweep). We report this explicitly as a \emph{pilot}, not a full
experiment it is a single, unreplicated data point we did not run the
Kaggle$\to$HandPD direction, additional seeds, or a weight sweep, and a
proper evaluation of domain-alignment methods for this problem is a
substantial undertaking in its own right that we leave to dedicated future
work. Gulrajani and Lopez-Paz's large-scale audit of domain-generalisation
algorithms is a direct motivation for treating this single run as a pilot
rather than a conclusive test under their controlled \texttt{DomainBed}
comparison, no domain-generalisation algorithm they evaluated (including
correlation-alignment and domain-adversarial methods) reliably outperforms
plain empirical risk minimisation by more than about one point once model
selection is held fixed across methods, which means a single-seed,
single-hyperparameter CORAL run of the kind we report here is exactly the
sort of result their audit warns is too noisy to interpret as a genuine
algorithmic effect one way or the other \cite{gulrajani2021}. Within that limited scope, the pilot is at minimum consistent with
the rest of this study's pattern none of the interventions tested so far
(preprocessing-based suppression or a single representation-alignment
pilot) close the honest-protocol gap on this data.

\begin{figure}[H]
\centering
\includegraphics[width=0.62\linewidth]{intervention_comparison_fig.png}
\caption{All interventions tested for HandPD$\to$Kaggle cross-cohort
transfer, EfficientNet-B0, four verified seeds (mean $\pm$ SD). Neither
background suppression nor the Deep CORAL pilot (a single unreplicated
run, not a full experiment) improves on the plain baseline for
this direction.}
\label{fig:interventions}
\end{figure}

\subsection{Recommendations for evaluating proposed de-leaking interventions}
Generalising from what this study needed in order to distinguish a genuine
null (Fork B: the intervention worked but did not help) from a failed
manipulation (Fork C: the intervention did not do what it was supposed to),
we suggest the following minimum protocol for any future work proposing a
preprocessing-based fix for identity or acquisition leakage in medical
image classification:
\begin{enumerate}[leftmargin=1.4em,itemsep=2pt]
  \item \textbf{Localise before you build.} Run a cheap, frozen-feature
        signal-localisation probe (identity re-identification and any
        known confound's decodability, under the masking modes relevant to
        the proposed intervention) before building the intervention,
        to establish whether it has a priori headroom.
  \item \textbf{Verify the manipulation, not just the outcome.} Re-run the
        same probe on the intervened data. A null outcome without a
        confirmed drop in the targeted signal is uninterpretable a
        positive outcome without a confirmed drop in the targeted signal
        may be a different, unidentified leak (Fork C).
  \item \textbf{Report the causal statistic with paired, subject-level
        uncertainty}, not independent point estimates for baseline and
        treatment, and evaluate in both the within-cohort and (where
        possible) cross-cohort external-validation regimes.
  \item \textbf{Rule out preprocessing artefacts} that could masquerade as
        the effect of interest (here, unequal resolution standing in for
        an acquisition marker).
  \item \textbf{Report negative and mixed results in full}, including
        directional but non-significant costs, rather than rounding to a
        clean null or omitting an intervention that did not work.
\end{enumerate}

\subsection{Limitations}
\label{sec:limitations}
First, our probe uses frozen ImageNet features as a proxy for what a
trained classifier could exploit because these features are not tuned to
handwriting, trace-intrinsic identity signal measured here is a floor, not
a ceiling, so the true trace-borne leak may be larger than reported.
Second, the trace/background split reuses a single first-pass segmentation
mask on HandPD the printed guide partially overlaps the pen trace, so the
trace-only view is not a perfectly clean separation, and we treat small
(\textless0.05) between-mode differences as within the noise floor of the
method rather than as evidence. Third, both cohorts are small (HandPD 18
healthy subjects Kaggle 28 subjects total), so single-comparison
confidence intervals are wide we mitigate this with subject-level
bootstrap statistics and by testing the effect across two independent
evaluation regimes (within-cohort and cross-cohort) rather than relying on
either alone. Fourth, we tested one suppression method (adaptive
thresholding after illumination flattening) a different background-removal
technique could in principle behave differently, though the identity-probe
result (Table~\ref{tab:probe-raw}) suggests any purely background-confined
intervention would face the same ceiling, since identity is not
background-confined to begin with. Fifth, the causal comparison
(Table~\ref{tab:causal}, Figure~\ref{fig:forest})
is restricted to EfficientNet-B0 we anchor on this backbone because our
own baseline replication (Table~\ref{tab:baseline}) found the other two
architectures too high-variance for a
two-condition paired comparison at this sample size, but this means the
result should be read as evidence for one stable architecture, not as a
claim already generalised across backbones. We do test both cross-cohort
transfer directions (HandPD$\rightarrow$Kaggle and Kaggle$\rightarrow$HandPD),
and both show the same wide, near-zero pattern. Sixth, cross-cohort
external validation is necessarily restricted to the spiral task, since it
is the only task name common to both cohorts we mitigate this by
replicating the within-cohort causal comparison on HandPD's and Kaggle's
other tasks (meander and wave respectively, Section~\ref{sec:othertasks}),
which show the same pattern, but a genuine cross-cohort test on a
non-spiral task is not possible with these two datasets. Seventh, HandPD
and Kaggle differ in class prevalence ($80.4\%$ vs.\ $53.6\%$ PD,
Section~\ref{sec:baseline}) we use AUC as the primary endpoint because it
is largely prevalence-invariant, but we did not run a prevalence-matched
subsampling sensitivity analysis to quantify how much of the cross-cohort
AUC gap, if any, is attributable to the prevalence mismatch itself rather
than to acquisition/domain shift this would be a natural addition to a
follow-up study. Relatedly, this study's cohort and trace-decodability
probes (Sections~\ref{sec:cohort-marker}--\ref{sec:rendering-invariance})
establish that geometric, trace-intrinsic signal is \emph{sufficient} to
decode cohort identity they do not, on their own, establish that this
signal is the cause of the cross-cohort AUC drop. HandPD and Kaggle
differ not only in acquisition pipeline but in population age, disease
stage, and drawing instructions, none of which this study measured or
controlled any of these could independently contribute to the honest
cross-cohort AUC drop through the classifier's disease-relevant features,
not only through the trace-intrinsic identity/acquisition fingerprint we
localise here. We show that background suppression and one domain-alignment
method do not close the gap, which is evidence against those two specific
mechanisms, but it is not evidence that the trace-intrinsic signal we
identify is the gap's sole or even primary cause; a designed study
(matched age/stage subsamples, or a prospective single-pipeline, multi-site
cohort) would be needed to adjudicate between these explanations, and we
report our causal language throughout as pertaining to the tested
interventions' failure to close the gap, not to a fully identified causal
mechanism for the gap itself. Eighth, the causal-lane bootstrap (Algorithm~\ref{alg:boot})
pools resampled deltas across seeds by concatenation, which mixes
subject-sampling and seed-to-seed variance without separating them a
nested or hierarchical bootstrap, or a mixed-effects model with seed as a
random effect, would give a more principled variance decomposition, and we
did not implement one here. Ninth, and most consequentially during
revision we discovered that the original evaluation harness run backing
this study's results had, for at least the HandPD spiral baseline,
executed with an under-specified fold count (3 folds rather than the
intended 10) relative to the harness's own documented defaults. Rather
than leave this an open question, we independently regenerated every
EfficientNet-B0 cell feeding Tables~\ref{tab:baseline} and~\ref{tab:causal}
under four freshly verified, matched seeds (999, 1001-1003) and explicit,
logged configuration ($k{=}10$, 40 epochs): the HandPD spiral within-cohort
baseline and causal comparison, and both raw and suppressed predictions
for both cross-cohort directions (HandPD$\to$Kaggle and Kaggle$\to$HandPD).
Every verified number confirms the same qualitative conclusions as the
original, under-specified run the within-cohort causal comparison remains
a Fork-B null with a smaller point estimate and a cleaner verdict
($P(\Delta\text{AUC}>0)$ moved from a borderline $0.054$ to an unambiguous
$0.121$); the HandPD$\to$Kaggle direction moved only modestly, both for
the raw cell ($0.685\to0.703$) and the paired causal comparison
($-0.009\to-0.040$, still CI-spans-zero) and the Kaggle$\to$HandPD
direction, the noisiest of the three by far (only 28 Kaggle subjects total
on either the training or test side depending on direction), moved further
in absolute point-estimate terms ($0.703\to0.647$ raw,
$+0.022\to+0.017$ paired) but the verified estimates are consistently
tighter than the originals (raw SD $0.130\to0.056$ paired CI actually
widened to $-0.404$ to $+0.309$, if anything a more honest reflection of
how underpowered this direction is) and remain well within the original
estimates' own wide uncertainty. No verified correction changed this
study's qualitative conclusions all three causal comparisons remain
Fork-B nulls only their precision and provenance. We disclose every
correction explicitly rather than silently revising the affected numbers.
The only cells of Table~\ref{tab:baseline} still carried over from the
original, unverified harness invocation are the VGG-16 and ResNet-50
cross-cohort cells (both directions) we did not extend the verification to
these secondary backbones given the one-job-per-day compute budget this
revision operated under and because the causal comparison in
Section~\ref{sec:stats} anchors on EfficientNet-B0 alone. Independently
regenerating these remaining cells would be a natural, lower-priority
follow-up, well behind the extensions discussed in
Section~\ref{sec:future}. Tenth, the positive-control marker experiment
(Section~\ref{sec:positive-control}) surfaced its own honest caveat worth
flagging on its own terms. Suppression collapsed the injected marker's
background-level decodability completely (AUC $0.991\to0.461$, i.e.\ to
chance), yet honest-protocol AUC recovered only part of the marker's
inflation ($0.994\to0.972$, against a marker-free baseline of $0.944$)
roughly $44\%$ of the injected effect, not all of it. We do not know
whether this residual reflects a genuine mismatch between what a
frozen-feature linear probe and a fine-tuned CNN are each sensitive to, a
subtle artefact introduced by the suppression reconstruction itself
distinct from the marker, or simply the five-seed design's own sampling
noise (we did not run a paired bootstrap CI for this comparison, unlike the
main causal-lane results, so we cannot rule out that the residual is not
statistically distinguishable from zero). Whichever explanation is
correct, this means a clean diagnostic-lane result does not guarantee the
causal lane's AUC movement will track it one-for-one in magnitude, which
tempers how much weight the diagnostic lane's clean collapse in the real,
non-marker suppression result (Table~\ref{tab:probe-supp}) can carry on
its own toward interpreting the corresponding Fork-B null
(Table~\ref{tab:causal}); it does not, however, change the sign or
qualitative interpretation of either result. Eleventh: as documented in
Section~\ref{sec:training}, the within-cohort P1/P2/LOSO protocols use
each fold's own held-out test split as the early-stopping validation set,
with no further inner validation split the reported checkpoint for
each fold is the epoch with the best AUC on that fold's own test data.
This is a mild form of model selection on the test set and can inflate
within-cohort AUC estimates somewhat, since the stopping point (though not
the training data itself) is chosen with knowledge of test-fold
performance. The cross-cohort comparisons this study's main causal
evidence are not affected, since their validation slice is carved from
the training cohort only. A cleaner design would nest a
separate inner validation split inside each within-cohort training fold
we did not implement this for the original baseline harness and disclose
it here as a limitation on every within-cohort number (P1, P2, LOSO)
reported in this study rather than silently correcting it, since doing so
would require retraining every within-cohort cell under a new splitting
scheme and was out of scope for this revision.

% ======================================================================
\section{Conclusion}
Subject identity and cross-cohort acquisition signal in this handwriting
data are not confined to the image background or template a
mechanistically verified background/template suppression removes the
background-specific component of both signals but leaves the honest
generalisation gap unchanged, because a comparable signal survives in the
drawn trace itself. This localises, rather than merely restates, the
leakage problem this study's own baseline replication establishes, and
rules out preprocessing-only masking as a sufficient fix. A single,
preliminary test of one domain-invariant training method (Deep CORAL) also
showed no improvement over the plain baseline, though this was only one
direction, one hyperparameter setting, and one unreplicated run
(Section~\ref{sec:future}), so we do not treat it as ruling out that
family of methods the way our suppression result rules out
preprocessing-only masking. We recommend that future
handwriting-based PD screening work pair any proposed de-leaking
intervention with a direct signal-localisation probe, not only an
accuracy comparison, and we point future mechanism work toward a more
thorough evaluation of domain-invariant training (broader hyperparameter
and direction coverage than we could run here) and multimodal (image +
dynamic signal) fusion as the next candidates.

% ======================================================================
\section*{Ethics statement}
This study performed a secondary analysis of two pre-existing, publicly
available, de-identified handwriting datasets (HandPD and the Kaggle
``Parkinson's Drawings'' set) no new human-subject data were collected.
Both datasets were collected and released by their original authors under
their own institutional approvals \cite{pereira2015,zham2017}. No additional
ethical approval was required for this secondary analysis.

\section*{Data and code availability}
The datasets are publicly available (HandPD from the Botucatu/UNESP dataset
page \cite{pereira2015}; the Kaggle ``Parkinson's Drawings'' set from Kaggle
\cite{zham2017}). All code implementing the baseline replication, the
masking probe, the suppression pipeline, and the paired statistical
comparison, together with the result data behind every table and figure and
a step-by-step reproduction guide, is released at
\url{https://github.com/jesuarockiaraj/localising-subject-identity-and-acquisition-leakage-in-Parkinson-handwriting-classifiers}.
A runnable demonstration
(\texttt{train\_demo\_checkpoint.py}, \texttt{demo\_predict.py}) trains one
EfficientNet-B0 classifier under the honest subject-grouped protocol and
reports its PD-probability output on full / background-only / trace-only
versions of a chosen image, so the trace-intrinsic result in
Table~\ref{tab:probe-supp} can be observed directly on individual
predictions rather than only in aggregate statistics.
\texttt{make\_worked\_example.py} reproduces Figure~\ref{fig:worked} from
the same checkpoint. \texttt{rendering\_invariance\_probe.py} reproduces
Tables~\ref{tab:rendering}--\ref{tab:rendering-supp}
(Sections~\ref{sec:rendering-invariance}--\ref{sec:rendering-invariance2}),
and \texttt{stratified\_reid.py} reproduces Table~\ref{tab:reid-strat} and
the label-shuffle control of Section~\ref{sec:reid-controls}.
\texttt{aspect\_ratio\_control.py} reproduces the pad-to-square check in
Section~\ref{sec:resolution-artefact}, and \texttt{handpd\_clean\_split.py}
reproduces the colour-separated pen-only result in
Section~\ref{sec:handpd-clean}. \texttt{verify\_guide\_color.py} reproduces
the cross-subject mask IoU check that establishes which colour is the
printed guide. \texttt{rendering\_invariance\_penonly.py} and
\texttt{rendering\_invariance\_penonly\_suppressed.py} reproduce the
guide-excluded rendering-invariance ladder (Table~\ref{tab:rendering-penonly})
on raw and suppressed images respectively, and
\texttt{background\_disease\_probe.py} and \texttt{compute\_suppressed\_p1.py}
reproduce the background-only disease-decodability and suppressed P1/P2
results discussed in Sections~\ref{sec:handpd-clean}--\ref{sec:limitations}.
Training configuration (Section~\ref{sec:training})
and package versions are fixed in the released code so that all reported
numbers can be regenerated end to end.

\section*{CRediT authorship contribution statement}
\textbf{Jesu Arockiaraj J}: Conceptualization, Methodology, Software,
Investigation, Formal analysis, Data curation, Visualization, Writing
original draft. 
\textbf{Arulkumar V}: Supervision, Methodology,
Validation, Writing  review \& editing. Both authors approved the final
version.

\section*{Declaration of competing interest}
On behalf of all
authors, the corresponding author states that there is no conflict of interest.
\section*{Funding}
This research received no external funding. Computational resources were
provided by Vellore Institute of Technology, Vellore, India.

\section*{Research involving human and/or animal participants}
Not applicable: this study performed a secondary analysis of pre-existing,
publicly available, de-identified data (see Ethics statement above) no
new human or animal participants were involved.

\section*{Informed consent}
Not applicable: no new participants were recruited, so no new informed
consent was obtained by the authors consent for the original data
collection was obtained by each dataset's original authors under their
own institutional approvals \cite{pereira2015,zham2017}.

\section*{Declaration of generative AI and AI-assisted technologies in the manuscript preparation process}
During the preparation of this work, the author(s) used AI Tools
to assist with data analysis scripting, figure generation, and drafting and
editing of the manuscript text. After using this tool, the author(s)
reviewed and edited the content as needed and take full responsibility for
the content of the published article.

% ======================================================================
\begin{thebibliography}{9}

\bibitem{chaibubneto2019}
E. Chaibub Neto et al., ``Detecting the impact of subject characteristics
on machine learning-based diagnostic applications,'' \emph{npj Digital
Medicine}, vol.\ 2, p.\ 99, 2019.

\bibitem{tampu2022}
I. E. Tampu, A. Eklund, N. Haj-Hosseini, ``Inflation of test accuracy due
to data leakage in deep learning-based classification of OCT images,''
\emph{Scientific Data}, vol.\ 9, p.\ 580, 2022.

\bibitem{gazda2022}
M. Gazda, M. Hire\v{s}, P. Drot\'{a}r, ``Multiple-fine-tuned convolutional
neural networks for Parkinson's disease diagnosis from offline
handwriting,'' \emph{IEEE Trans.\ Systems, Man, and Cybernetics: Systems},
vol.\ 52, no.\ 1, pp.\ 78--89, 2022.

\bibitem{kamran2021}
I. Kamran, S. Naz, I. Razzak, M. Imran, ``Handwriting dynamics assessment
using deep neural network for early identification of Parkinson's
disease,'' \emph{Future Generation Computer Systems}, vol.\ 117,
pp.\ 234--244, 2021.

\bibitem{pereira2015}
C. R. Pereira, D. R. Pereira, F. A. Silva, C. Hook, S. A. T. Weber,
L. A. M. Pereira, J. P. Papa, ``A step towards the automated diagnosis of
Parkinson's disease: Analyzing handwriting movements,'' in \emph{Proc.\
IEEE 28th Int.\ Symp.\ Computer-Based Medical Systems (CBMS)}, S\~{a}o
Carlos, Brazil, 2015, pp.\ 171--176; HandPD dataset,
\url{https://wwwp.fc.unesp.br/~papa/pub/datasets/Handpd/}.

\bibitem{zham2017}
P. Zham, D. K. Kumar, P. Dabnichki, S. Poosapadi Arjunan, S. Raghav,
``Distinguishing different stages of Parkinson's disease using a
composite index of speed and pen-pressure of sketching a spiral,''
\emph{Frontiers in Neurology}, vol.\ 8, p.\ 435, 2017. Kaggle
``Parkinson's Drawings'': \url{https://www.kaggle.com/datasets/kmader/parkinsons-drawings}.

\bibitem{brookshire2024}
G. Brookshire, J. Kasper, N. M. Blauch, Y. C. Wu, R. Glatt, D. A. Merrill,
S. Gerrol, K. J. Yoder, C. Quirk, C. Lucero, ``Data leakage in deep
learning studies of translational EEG,'' \emph{Frontiers in Neuroscience},
vol.\ 18, 1373515, 2024.

\bibitem{rosenblatt2024}
M. Rosenblatt, L. Tejavibulya, R. Jiang, S. Noble, D. Scheinost, ``Data
leakage inflates prediction performance in connectome-based machine
learning models,'' \emph{Nature Communications}, vol.\ 15, no.\ 1,
p.\ 1829, 2024.

\bibitem{wen2020}
J. Wen, E. Thibeau-Sutre, M. Diaz-Melo, J. Samper-Gonz\'{a}lez,
A. Routier, S. Bottani, D. Dormont, S. Durrleman, N. Burgos, O. Colliot,
``Convolutional neural networks for classification of Alzheimer's
disease: Overview and reproducible evaluation,'' \emph{Medical
Image Analysis}, vol.\ 63, 101694, 2020.

\bibitem{alzheimerscoping2025}
``Data leakage in deep learning for Alzheimer's disease diagnosis: a
scoping review of methodological rigor and performance inflation,''
\emph{Diagnostics}, vol.\ 15, 2348, 2025.

\bibitem{geirhos2020}
R. Geirhos, J.-H. Jacobsen, C. Michaelis, R. Zemel, W. Brendel, M. Bethge,
F. A. Wichmann, ``Shortcut learning in deep neural networks,'' \emph{Nature
Machine Intelligence}, vol.\ 2, pp.\ 665--673, 2020.

\bibitem{degrave2021}
A. J. DeGrave, J. D. Janizek, S.-I. Lee, ``AI for radiographic COVID-19
detection selects shortcuts over signal,'' \emph{Nature Machine
Intelligence}, vol.\ 3, pp.\ 610--619, 2021.

\bibitem{zech2018}
J. R. Zech, M. A. Badgeley, M. Liu, A. B. Costa, J. J. Titano,
E. K. Oermann, ``Variable generalization performance of a deep learning
model to detect pneumonia in chest radiographs: a cross-sectional study,''
\emph{PLOS Medicine}, vol.\ 15, no.\ 11, e1002683, 2018.

\bibitem{roberts2021}
M. Roberts, D. Driggs, M. Thorpe, et al., ``Common pitfalls and
recommendations for using machine learning to detect and prognosticate for
COVID-19 using chest radiographs and CT scans,'' \emph{Nature Machine
Intelligence}, vol.\ 3, pp.\ 199--217, 2021.

\bibitem{badgeley2019}
M. A. Badgeley, J. R. Zech, L. Oakden-Rayner, B. S. Glicksberg, M. Liu,
W. Gale, M. V. McConnell, B. Percha, T. M. Snyder, J. T. Dudley, ``Deep
learning predicts hip fracture using confounding patient and healthcare
variables,'' \emph{npj Digital Medicine}, vol.\ 2, p.\ 31, 2019.

\bibitem{kaufman2012}
S. Kaufman, S. Rosset, C. Perlich, O. Stitelman, ``Leakage in data mining:
Formulation, detection, and avoidance,'' \emph{ACM Transactions on
Knowledge Discovery from Data}, vol.\ 6, no.\ 4, pp.\ 1--21, 2012.

\bibitem{varoquaux2022}
G. Varoquaux, V. Cheplygina, ``Machine learning for medical imaging:
methodological failures and recommendations for the future,'' \emph{npj
Digital Medicine}, vol.\ 5, p.\ 48, 2022.

\bibitem{ganin2015}
Y. Ganin, V. Lempitsky, ``Unsupervised domain adaptation by
backpropagation,'' in \emph{Proc.\ ICML}, 2015.

\bibitem{sun2016coral}
B. Sun, K. Saenko, ``Deep CORAL: Correlation alignment for deep domain
adaptation,'' in \emph{Proc.\ ECCV Workshops}, 2016.

\bibitem{packhauser2022}
K. Packh\"auser, S. G\"undel, N. M\"unster, C. Syben, V. Christlein,
A. Maier, ``Deep learning-based patient re-identification is able to
exploit the biometric nature of medical chest X-ray data,''
\emph{Scientific Reports}, vol.\ 12, p.\ 14851, 2022.

\bibitem{pereira2016newhandpd}
C. R. Pereira, D. R. Pereira, F. A. Silva, J. P. Masieiro, S. A. T. Weber,
C. Hook, J. P. Papa, ``A new computer vision-based approach to aid the
diagnosis of Parkinson's disease,'' \emph{Computer Methods and Programs in
Biomedicine}, vol.\ 136, pp.\ 79--88, 2016.

\bibitem{drotar2016pahaw}
P. Drot\'{a}r, J. Mekyska, I. Rektorov\'{a}, L. Mas\'{a}rov\'{a},
Z. Sm\'{e}kal, M. Faundez-Zanuy, ``Evaluation of handwriting kinematics
and pressure for differential diagnosis of Parkinson's disease,''
\emph{Artificial Intelligence in Medicine}, vol.\ 67, pp.\ 39--46, 2016.

\bibitem{impedovo2019}
D. Impedovo, G. Pirlo, ``Dynamic handwriting analysis for the assessment
of neurodegenerative diseases: a pattern recognition perspective,''
\emph{IEEE Reviews in Biomedical Engineering}, vol.\ 12, pp.\ 209--220,
2019.

\bibitem{saeb2017}
S. Saeb, L. Lonini, A. Jayaraman, D. C. Mohr, K. P. Kording, ``The need
to approximate the use-case in clinical machine learning,''
\emph{GigaScience}, vol.\ 6, no.\ 5, gix019, 2017.

\bibitem{kapoor2023}
S. Kapoor, A. Narayanan, ``Leakage and the reproducibility crisis in
machine-learning-based science,'' \emph{Patterns}, vol.\ 4, no.\ 9,
100804, 2023.

\bibitem{gulrajani2021}
I. Gulrajani, D. Lopez-Paz, ``In search of lost domain generalization,''
in \emph{Proc.\ ICLR}, 2021.

\bibitem{lakens2017}
D. Lakens, ``Equivalence tests: a practical primer for $t$ tests,
correlations, and meta-analyses,'' \emph{Social Psychological and
Personality Science}, vol.\ 8, no.\ 4, pp.\ 355--362, 2017.

\end{thebibliography}

\end{document}
